\documentclass{article}
\setlength{\topmargin}{-0.25 in}
\setlength{\textwidth}{6.5in}
\setlength{\textheight}{8.5 in}
\setlength{\oddsidemargin}{0in}
\setlength{\evensidemargin}{0in}
\setlength{\parindent}{0.0in}
\setlength{\parskip}{7pt plus 1pt}

\usepackage{textcomp}
\usepackage{bcprules}
\usepackage{xcolor}				% Allow for color (annotations)
\usepackage{stmaryrd}
\usepackage{amssymb}
\usepackage{hyperref}

\def\sectionautorefname{Section}
\def\subsectionautorefname{Section}
\def\subsubsectionautorefname{Section}
\def\tableautorefname{Table}
\def\figureautorefname{Figure}

\def\tdm#1{\textcolor{blue}{\textsf{TDM: #1}}}
\def\ab#1{\textcolor{red}{\textsf{AB: #1}}}
\def\tf#1{\textcolor{purple}{\textsf{TF: #1}}}

\newcommand{\maybe}{{\bf\footnotesize{?}}\kern-1pt}

\newcommand{\eval}[3]{{\tt #1} $\vdash$ {\tt #2} $\Downarrow$ {\tt #3}}
\newcommand{\stay}[2]{stay({\tt #1}) $=$ {\tt #2}}
\newcommand{\istay}[2]{i-stay({\tt #1}) $=$ {\tt #2}}
\newcommand{\toSolver}[2]{{\tt E} $\vdash$ {\tt #1} $\leadsto$ $#2$}
\newcommand{\fromSolver}[2]{$#1$ $\leadsto$ {\tt #2}}
\newcommand{\csolve}[2]{$#1$ $\models$ $#2$}
\newcommand{\solve}[2]{{\tt #1} $\models$ {\tt #2}}
\newcommand{\cfail}[1]{$\not\models$ $#1$}
\newcommand{\fail}[1]{$\not\models$ {\tt #1}}
\newcommand{\step}[5]{{\tt <#1|#2|#3>} $\longrightarrow$
  {\tt <#4|#5>}}
\newcommand{\tc}[3]{{\tt #1} $\vdash$ {\tt #2} : {\tt #3}}
\newcommand{\tcC}[2]{{\tt #1} $\vdash$ {\tt #2}}

\newcommand{\evalH}[4]{{\tt #1;#2} $\vdash$ {\tt #3} $\Downarrow$ {\tt #4}}
\newcommand{\solveH}[3]{{\tt #1;#2} $\models$ {\tt #3}}
% \newcommand{\solveH}[3]{{\tt #1;#2} $\models_{E,H}$ {\tt #3}}
% \newcommand{\solveH}[3]{{\tt #1} $\models_{E,H}$ {\tt #3}}
\newcommand{\failH}[1]{$\not\models_{E,H}$ {\tt #1}}
\newcommand{\stepH}[9]{{\tt <#1|#2|#3|#4|#5>} $\longrightarrow$
  {\tt <#6|#7|#8|#9>}}
\newcommand{\tcH}[4]{{\tt #1;#2} $\vdash$ {\tt #3} : {\tt #4}}
\newcommand{\tcHC}[3]{{\tt #1;#2} $\vdash$ {\tt #3}}
\newcommand{\toSolverH}[2]{{\tt E,H} $\vdash$ {\tt #1} $\leadsto$ $#2$}
\newcommand{\fromSolverH}[3]{$#1$ $\leadsto$ {\tt #2,#3}}

\newcommand{\toSolverOO}[6]{{\tt #1,#2} $\vdash$ {\tt #3} $\leadsto$ {\tt #4,#5} $\mid$ $#6$}

\newcommand{\inferrule}[3]{\infrule[#1]{\mbox{#2}}{\mbox{#3}}}
\newcommand{\inferax}[2]{\infax[#1]{\mbox{#2}}}

\newcommand{\judgment}[1]{\begin{flushleft} \fbox{#1}\end{flushleft}}

\newcommand{\receq}{\stackrel{{\tt r}}{{\tt =}}}

\newcommand{\lookup}[3]{\textit{lookup}({\tt #1},{\tt #2}) = ({\tt x$_1$} $\cdots$ {\tt x$_n$},{\tt #3})}


\newcommand{\inline}[9]{{\tt <#1,#2,#3,#4,#5,#6>} $\leadsto$ {\tt <#7,#8,#9>}}


\newcommand\enter[9]{%
\def\argA{#1}\def\argB{#2}\def\argC{#3}\def\argD{#4}\def\argE{#5}%
\def\argF{#6}\def\argG{#7}\def\argH{#8}\def\argI{#9}%
\entercontinued
}
\newcommand\entercontinued[2]{%
  \textit{enter}({\tt \argA},{\tt \argB},{\tt \argC},{\tt \argD},{\tt
    \argE},{\tt \argF},{\tt x$_1$} $\cdots$ {\tt x$_n$},{\tt e$_1$} $\cdots$
  {\tt e$_n$}) = ({\tt \argG},{\tt \argH},{\tt \argI},{\tt #1},{\tt #2})
}

\newcommand\stepHO[9]{%
\def\argA{#1}\def\argB{#2}\def\argC{#3}\def\argD{#4}\def\argE{#5}%
\def\argF{#6}\def\argG{#7}\def\argH{#8}\def\argI{#9}%
\stepHOcontinued
}
\newcommand{\stepHOcontinued}[2]{%
  {\tt <\argA|\argB|\argC|\argD|\argE|\argF>} $\longrightarrow$ {\tt <\argG|\argH|\argI|#1|#2>}}


\newcommand\evalBig[9]{%
\def\argA{#1}\def\argB{#2}\def\argC{#3}\def\argD{#4}\def\argE{#5}%
\def\argF{#6}\def\argG{#7}\def\argH{#8}\def\argI{#9}%
\evalBigcontinued
}
\newcommand{\evalBigcontinued}[2]{%
  {\tt <\argA|\argB|\argC|\argD|\argE|\argF>} $\Downarrow$ {\tt <\argG|\argH|\argI|#1|#2>}}

\newcommand{\resolve}[4]{{\tt #3;#4} $\models_{\tt #1}$ {\tt #2}}

\newcommand{\E}{$\mathbb{E}$}
\newcommand{\C}{$\mathbb{C}$}
\newcommand{\I}{$\mathbb{I}$}
\renewcommand{\H}{$\mathbb{H}$}
\newcommand{\reinline}[6]{{\tt <#1,#2,#3,#4>} $\leadsto$ {\tt <#5,#6>}}

\sloppy

\begin{document}

\title{Developing a Formal Semantics for Babelsberg: \\ 
  A Step-by-Step Approach}
\author{Tim Felgentreff, Todd Millstein, and Alan Borning}
\date{revised version of VPRI Technical Report TR-2014-002 \\ 
draft of \today}

% \date{VPRI Technical Report TR-2014-002 \\ September 2014}

% changed title to make it different from the title of a possible
% later conference submission

\maketitle

\section{Introduction}

Babelsberg \cite{felgentreff-babelsberg-tr} is a family of object
constraint languages, with current instances being Babelsberg/R (a
Ruby extension), Babelsberg/JS \cite{felgentreff-ecoop-2014} (a
Javascript extension), and Babelsberg/S (a Squeak extension). The
Babelsberg design has evolved alongside its implementations, driven in
part by practical considerations about the expectations of programmers
familiar with the underlying host language.  This fact, along with the
complexities of integrating objects, state, and constraints, have led
to a number of the semantic choices being muddled with implementation
specifics.  There have also been a number of long-standing, confusing
issues with respect to constraints and object identity, how to
represent assignment, and the appropriate restrictions on expressions
that define constraints.  In an effort to understand these better and
to provide a complete design that instances of Babelsberg can
implement, we give here a formal operational semantics for Babelsberg.

We've found it helpful to approach the problem incrementally, first
devising a formal semantics for a very simple constraint imperative
language, and then building up to a full object constraint language.
In this memo we present the semantics in that fashion as well.  The
languages are as follows:

\begin{itemize}

\item Babelsberg/PrimitiveTypes (the only datatypes are booleans plus
  a set of primitive types, such as integers, reals, and strings).
  The concrete variant we use in informal examples is Babelsberg/Reals.
\item Babelsberg/Records (Babelsberg with immutable records, along
  with primitive types).
\item Babelsberg/UID (Babelsberg with mutable records that live on the
  heap and so have an identity, can be aliased, etc., as well as
  primitive types).
\item Babelsberg/Objects (Babelsberg with mutable objects, classes,
  methods, messages, inheritance, and object-oriented constraint
  definitions).
\end{itemize}

In each case, we first provide an informal discussion and examples,
and then the formal semantics.  The primary audience for this
admittedly long tech report is first ourselves --- we have clarified
many aspects of the language as a result of working out the formalism
--- and second, other researchers who are interested in the details.
Subsequently we plan to write a standard-length paper for a
programming language conference that should be of more general
interest, based on this material.  (This subsequent paper will focus just
on Babelsberg/Objects, perhaps with a brief introductory example using
Babelsberg/Reals, and will also include proofs of some of the
interesting properties of the language.)

\section{Motivation}

Our formal semantics is intended to provide a complete semantics of
Babelsberg that can be used to inform practical implementations of
the language.  It is meant to be as simple as possible, while still 
encompassing the
major design decisions needed to guide language implementers. Because
Babelsberg is a design to provide object-constraint programming in an
object-oriented host language, the semantics omits some constructs
such as exception handling for constraint solver failures and
syntactic sugar that are intended to be inherited from the host
language. The semantics instead focuses on the expression of standard
object-oriented constructs that need to be modified to support the
Babelsberg design.

An overarching design goal is that in the absence of constraints,
Babelsberg should be a standard object-oriented language.  We
have thus resisted the temptation to add other interesting features;
instead we aim to make the smallest possible changes while still accommodating
constraints in a clean and powerful way.

We require that our design support a useful and expressive language
for constraints, integrated with the host object oriented language in
a clean way.  Beyond this, however, a heuristic for the design is that
when there is a choice, we favor simplicity over power (at least power
to do interesting operations in the language that have nonetheless not
yet proven themselves useful in practice, in particular with respect
to object identity and type).  The constraint imperative language
Kaleidoscope \cite{freeman-benson-oopsla-1990, freeman-benson-phd,
  freeman-benson-iccl-1992, freeman-benson-ecoop-1992}, particularly
the early versions, was arguably too complex in part because it was
too powerful in interesting ways of just this kind, making it
difficult to understand what the result of a program might be and also
difficult to implement efficiently.  Trying to follow this heuristic
also makes the design more independent of the particular host language
and solvers used in implementing it, because only a small set of basic
operations have to be adapted.

Some significant clarifications and simplifications of Babelsberg that
have come out of this work on formalizing the language are:

\begin{itemize}
\item A clearer understanding of the interactions among constraints
  on identity, types, and values.
\item The addition of structural compatibility checks to tame the
  power of the solver with respect to changing object structure and
  type to satisfy constraints.
\item The addition of \emph{value classes}.  (Instances of value
  classes are immutable objects for which object identity is not
  significant.)  Value classes play a key role in giving a clear specification
  of the requirements for expressions that define constraints (in
  particular that they must be side effect free), while still supporting
  useful programming idioms.
\item A set of restrictions on identity constraints that make it
  easier to reason about object identity and type.  In particular, any
  change to the identity of the object referred to by a variable flows
  clearly from a single assignment statement, and is deterministic ---
  there are never multiple correct solutions to the identity
  constraints.  This also implies that any change to the type of a
  variable must similarly flow from a single assignment statement and
  be deterministic.
\end{itemize}

\section{Constraints}
\label{sec:constraints}

The semantics for each of these languages includes a step in which we
assert that some set of values for variables is a correct solution to
a collection of hard and soft constraints.  There is a constraint
solver that is a black box as far as the rest of the formal semantics
is concerned, and that handles all the issues regarding finding a
solution, dealing with conflicting soft constraints, and so forth.
The solver should be sound but may be incomplete (i.e., it should
never return an incorrect answer, but might respond that the set of
constraints is too difficult for it to determine whether or not there
is a solution).

We use the semantics for hard and soft constraints presented in
\cite{borning-lisp-symbolic-computation-1992}.  An earlier paper on
Babelsberg \cite{felgentreff-babelsberg-tr} has a description of the
relevant theory as well.  The knowledge about how to handle hard
and soft constraints is left entirely up to the solver and doesn't
enter into the formal semantics, so the discussion in the rest of this
section is intended to help the reader understand the informal
examples --- it isn't relevant to the formal semantics.

The solver should find a single best solution --- if there are
multiple solutions, the solver is free to pick any one of them.
(Providing answers rather than solutions, i.e., results such as $10
\leq x \leq 20$ rather than a single value for $x$, and backtracking
among multiple answers, as available in for example constraint logic
programming \cite{jaffar-popl-1987}, is left for future work ---
see Appendix \ref{sec:adding-solvers}.)

The way we trade off conflicting soft constraints is defined by a
\emph{comparator}.  The solver encodes the comparator being used,
making it irrelevant to the formal semantics --- however, in
presenting the informal program examples we'll sometimes need to
specify which is used, and for this reason we include a brief
discussion of comparators.

The two most relevant comparators are
locally-predicate-better (LPB) and weighted-sum-better (WSB)\@.
Locally-predicate-better only cares whether a constraint is satisfied
or not, not how far off the value is from the desired one.  Any
Pareto-optimal solution is acceptable.  For example, a solution that
satisfies one weak constraint A and violates three weak constraints B,
C, and D is OK, as long as there isn't a solution that satisfies both
A and some additional constraint, even if there is another solution
that satisfies B, C, and D but not A\@.  The DeltaBlue
solver~\cite{freeman-benson-cacm-1990} finds LPB solutions.
Weighted-sum-better considers the weighted sum of the errors of
constraints with a given soft priority, and picks a solution that
minimizes the sum.  If there is more than one solution and there are
additional lower priority constraints, we then consider the
lower-priority ones to winnow down the possible solutions, priority by
priority.  For this comparator we need an error in satisfying a
constraint, which should be 0 iff the constraint is satisfied.
Cassowary \cite{badros-tochi-2001} finds WSB solutions.

Here are two examples.  (These are described from the point of view of
the declarative theory of hard and soft constraints, not with respect
to how an actual solver can find that solution.)

\begin{tabular}{rl}
required & $x+y=10$ \\
strong   & $x=8$ \\
weak     & $y=0$
\end{tabular}

The required constraint has an infinite number of solutions.  When we
winnow these down to solutions that satisfy the strong constraint,
there is only one left: $x=8$, $y=2$.  This is both a LPB and a WSB
solution.  The weak constraint has no impact on the solution in this
case.

Now consider:

\begin{tabular}{rl}
required & $x+y=10$ \\
strong   & $x \geq 5$ \\
weak     & $y=20$
\end{tabular}

We'll only consider the WSB comparator this time, since it is more
suitable for use with inequalities.  (DeltaBlue does not handle inequalities.)
Again, the required constraint has an infinite number of solutions.
We winnow these down with the strong constraint to all $x \in
[5,\infty)$, $y$ such that $x+y=10$.   The weak constraint is
  unsatisfiable, so we minimize its error, resulting in the solution
  $x=5$, $y=5$.

Strict inequality constraints with metric comparators can be
problematic in the presence of soft constraints, since they can lead
to no solutions.  Consider:

\begin{tabular}{rl}
required & $x>10$ \\
weak     & $x=5$
\end{tabular}

This set of constraints has no solution --- for any potential solution
that satisfies $x>10$, we can find another that better satisfies
$x=5$.  For this reason, our examples usually use non-strict
inequalities. 

The ``required'' priority is special, in that those constraints must
be satisfied in any solution.  Both the semantics of hard and soft
constraints, and the Cassowary and DeltaBlue solvers, can handle
arbitrary numbers of soft priorities.  However, for simplicity in the
remainder of this note, we only use three, namely ``strong,''
``medium,'' and ``weak''.

We can also annotate variables used in constraints as
\emph{read-only}.  Intuitively, when choosing the best solutions to a
set of constraints with priorities, constraints should not be allowed
to affect the choice of values for their read-only variables, i.e.,
information can flow out of the read-only variables, but not into
them.  Read-only annotations provide an important tool in a practical
language for guiding the behavior of the constraint solver.  However,
they don't present any particular issues for the formal semantics ---
we would just pass them on to the constraint solver to handle.
Therefore, for simplicity we've dropped the extra rules to handle them
(since they were just copies of the rules for variables without the
read-only annotation).

\subsection{Conjunctions and Disjunctions of Constraints}

In general, a constraint might consist of conjunctions, disjunctions,
and negations of atomic constraints.  For a conjunction or
disjunction, if there is a priority, it applies to the entire
constraint, not to components.  Thus this is legal:

strong $(x=3 \vee x=4)$

but this is not allowed:

(strong $x=3$) $\vee$ (weak $x=4$)

Only some solvers, such as Z3, can accommodate disjunctions and
explicit conjunctions of constraints.  For DeltaBlue and Cassowary,
conjunctions of constraints are implicitly specified by feeding
multiple constraints to the solver, while disjunctions aren't allowed.
However, as noted above, the solver is a black box as far as the rest
of the formal semantics is concerned.

\subsection{Taming Identity Constraints}
\label{sec:taming-identity-constraints}

We introduce object identity in the Babelsberg/UID language, and
continue to use it in our final language Babelsberg/Objects.  A
central issue in the design of a constraint object language is the
interaction between constraints and object identity.  Our earlier
experience with Kaleidoscope suggests that it is all too easy to make
constraints on object identity and types be too powerful, so that they
lead to non-obvious consequences.  This makes programs more difficult
for the programmer to understand.  To tame the power of constraints on
object identity and type, we set the following goals for this aspect
of our design:

\begin{enumerate}
\item In those languages that include object identity, we want to
  support explicit but straightforward identity constraints of the
  form \verb|x==y| (i.e., that variables \verb|x| and \verb|y| refer
  to the same object).
\item\label{item:deterministic} The solution to the constraints should
  be deterministic as far as object identity and type are concerned
  --- there should never be multiple correct solutions in which a
  given variable refers to objects with different identities in the
  different solutions.
\item Any change to the identities of the objects referred to by
  variables should flow from an assignment statement --- the constraint
  solver should not otherwise alter object identities.
\item\label{item:solve-identity-then-values} To make Babelsberg
  programs more understandable for the programmer, we want to be able
  to reason about the constraints so that we can first solve all the
  identity constraints (with a deterministic solution, in keeping with
  Goal \ref{item:deterministic}), and then the value constraints.  A
  design that meets this goal of course benefits the language
  implementer as well.
\end{enumerate}

In our intermediate language Babelsberg/Records, in which we introduce
structured data in the form of immutable records but not yet object
identity, we want analogous properties to hold: the solution to the
constraints should be deterministic as far as object structure is
concerned, and any changes to object structure should flow from an
assignment statement.

%% This subsection commented out -- we don't need multiple comparators in
%% the current version, since we solve for identity separately from value
%% and don't intermix these constraints.
%% \subsection{Using Multiple Comparators}
%% \label{sec:multiple-comparators}

%% The semantics for hard and soft constraints used in work to date has
%% always employed a single comparator to specify how to trade off
%% conflicting soft constraints.  However, we can extend the theory to
%% allow multiple comparators --- different priorities can use different
%% comparators.  For example, we might use LPB for comparing the
%% solutions to strong constraints, and WSB for the weak ones.  We will
%% use this extension later so that we can mix constraints on object
%% identity and object values.  (Identity constraints will be first
%% introduced for Babelsberg/UID (\autoref{sec:babelsberg-uid}),
%% and also used for Babelsberg/Objects (\autoref{sec:babelsberg-objects}.)
%% An identity constraint is either
%% satisfied or it's not --- a WSB comparator isn't appropriate for them.
%% But there may also be constraints on values, for example on the $x$
%% and $y$ fields of points, for which we do want to use a WSB
%% comparator.  As long as we give different priorities for the identity
%% and the value constraints, this extension allows doing this.

\section{Babelsberg/Reals and Babelsberg/PrimitiveTypes}
\label{sec:babelsberg-primitivetypes}

We start with a very basic language, Babelsberg/Reals, that has only
primitive values. In Babelsberg, constraints are expressions that
return a boolean -- the constraint solver's task is to find values for
the variables in the constraint's expression such that it
evaluates to true. So Boolean is a required datatype for
all Babelsberg languages. In addition, we add reals as a second
primitive type.

The evaluation model for Babelsberg/Reals is mostly standard for
ordinary expressions and statements.  In contrast, a statement that
adds a constraint starts with a duration, namely \verb|always| or
\verb|once|.  The expression following the duration is taken
unevaluated and added to the \emph{constraint store}.  For
\verb|always|, the constraint remains for the duration of the
program's execution; for \verb|once|, it is removed after the solver
finds a solution to the current set of constraints.  The only wrinkle
in the evaluation model for ordinary statements is that for an
assignment statement, we evaluate the expression on the right hand
side of the assignment, constrain the variable on the left hand side
of the assignment to be equal to the result using a \verb|once|
constraint, and then turn all the constraints over to the solver.
Doing this ensures that assignment interacts correctly with other
constraints.

Here are some examples.

\begin{verbatim}
x := 3;
x := 4;
always x>=10
\end{verbatim}

%% Old version, with imperative execution mode and constraint
%% construction mode.  (We might need the ideas about closures and the
%% environment of definition later, for Babelsberg/Objects.)
%%
%% Babelsberg/Reals, as with the other Babelsberg languages, uses two different
%% execution modes: \emph{imperative execution mode} and \emph{constraint
%%   construction mode}.  For ordinary statements, expressions are
%% evaluated in imperative execution mode.  This is mostly standard,
%% except that assignment is syntactic sugar for a \texttt{once}
%% constraint (i.e., a constraint that is asserted and then retracted).
%% This \verb|once| constraint is formed by evaluating the expression on
%% the right hand side of the assignment, and then constraining the
%% variable on the left hand side of the assignment to be equal to the
%% result.  In contrast, constraint construction mode is used with
%% constraints that are created by an explicit \verb|always| or
%% \verb|once| constraint.  The semantic model includes a
%% \emph{constraint store} consisting of closures corresponding to
%% constraints encountered in an \verb|always| or \verb|once| statement:
%% when an \verb|always| or \verb|once| statement is encountered, the
%% interpeter (as specified by the formal semantics) constructs a closure
%% over the expression's environment of definition and saves it in the
%% constraint store.  For \verb|always|, the closure remains for the
%% duration of the program's execution; for \verb|once|, it is removed
%% after the constraints are satisfied.

After evaluating the first statement we hand the following \verb|once|
constraint to the solver to find a value for \texttt{x}:

\begin{tabular}{rl}
required & $x=3$
\end{tabular}

The solver finds a value for $x$, which is then used to update the
environment to be $\texttt{x} \mapsto \texttt{3}$.

Note that programs, as well as the variable names and values in the
environment, are written in fixed pitch font.  For contrast, we write
the constraints that are handed to the solver in math font.

Continuing with the example, after the second statement we hand the
following constraints to the solver:

\begin{tabular}{rl}
weak     & $x=3$ \\
required & $x=4$
\end{tabular}

The weak $x=3$ constraint is the stay constraint that \texttt{x}
retain its previous value, while the required $x=4$ constraint comes
from the second assignment.  This has the solution $x=4$, resulting in
a new environment $\texttt{x} \mapsto \texttt{4}$.  

The third statement adds the \texttt{always} constraint to the
constraint store.  So after that statement we have the following
constraints:

\begin{tabular}{rl}
weak     & $x=4$ \\
required & $x \geq 10$
\end{tabular}

If we use a metric comparator such as weighted-sum-better or
least-squares-better, we get the solution $x=10$, since this minimizes
the error for the weak constraint.  If we use
locally-predicate-better, then every $x \in [10,\infty)$ is a solution,
and the system is free to select any of them.  (However, as noted in
Section \ref{sec:constraints}, typically we wouldn't use LPB if we
have inequalities.)

The following example illustrates using the same variable on the left
and right hand sides of an assignment statement, as well as the
interaction of assignments with \verb|always| constraints.
\begin{verbatim}
x := 3;
y := 0;
always y = x+100;
x := x+2
\end{verbatim}

After evaluating the first two statements and solving the resulting
constraints, the environment has the binding 
$\texttt{x} \mapsto \texttt{3}$, $\texttt{y} \mapsto 0$.  The third
statement causes the constraint \verb|always y = x+100| to be added to
the constraint store.  We then hand the following constraints to the
solver to find values for \texttt{x} and \texttt{y}:

\begin{tabular}{rl}
weak     & $x=3$ \\
weak     & $y=0$ \\
required & $y=x+100$
\end{tabular}

This has multiple possible solutions, and the solver can select any
one of them.  Suppose it picks $\texttt{x} \mapsto 3$, 
$\texttt{y} \mapsto 103$.

After evaluating the next statement, we have the following constraints:

\begin{tabular}{rl}
weak     & $x=3$ \\
weak     & $y=103$ \\
required & $y=x+100$ \\
required & $x=5$
\end{tabular}

The first two constraints are the weak stays on \verb|x| and \verb|y|,
the third comes from the always constraint in the constraint store,
and the fourth comes from the assignment \verb|x:=x+2| (where we
evaluated \verb|x+2| in the old environment to get 5).  After solving
these constraints, we have $\texttt{x} \mapsto 5$, 
$\texttt{y} \mapsto 105$.

New variables must be created with an assignment statement.  Thus
the following program is illegal --- we would need to create
\verb|x| first with an assignment statement before adding the
\verb|always| constraint.

\begin{verbatim}
always x=10;
\end{verbatim}

Requiring that variables be created before equating them with
something can be annoying for the programmer.  We decided that indeed
the above program is illegal in the formal semantics.  However, in
practical implementations, we can have a shortcut to allow it.  It
only works for = constraints where one side is a new variable and
all variables on the other side already exist.  The behavior is
that the system creates the new variable, evaluates the other side
of the equality, assigns it to the new variable, and then adds the
equality constraint as an \verb|always| constraint.

An interesting aspect of this semantics, both as presented informally
above and in the formalism that follows, is that we no longer model
assignment as a constraint on variables at different times, as was
done in for example the first Babelsberg paper
\cite{felgentreff-babelsberg-tr}.  See Appendix
\ref{sec:perturbation-vs-refinement} for a discussion of these
alternatives.  (This is relegated to the appendix since, while
interesting, it is largely orthogonal to the task at hand of providing
a formal semantics for Babelsberg).

The program might also introduce simultaneous equations and
inequalities.  For example:
\begin{verbatim}
x := 0;
y := 0;
z := 0;
always x+y+2*z = 10;
always 2*x+y+z = 20;
x := 100
\end{verbatim}

Assuming the solver can solve simultaneous linear equations,
after the final assignment we will have 
$\texttt{x} \mapsto \texttt{100}$,
$\texttt{y} \mapsto \texttt{-270}$,
$\texttt{z} \mapsto \texttt{90}$.

As an example of unsatisfiable constraints, consider:
\begin{verbatim}
x := 5;
always x<=10;
x := x+15
\end{verbatim}

After evaluating the first statement the environment includes the
binding $\texttt{x} \mapsto \texttt{5}$.  After evaluating the
statement that generates the {\tt always} constraint, we solve the
constraints

\begin{tabular}{rl}
weak     & $x=5$ \\
required & $x \leq 10$
\end{tabular}

This has the solution $x=5$.  Then we evaluate the last
assignment, resulting the constraints

\begin{tabular}{rl}
weak     & $x=5$ \\
required & $x \leq 10$ \\
required & $x = 20$
\end{tabular}

Note that the required $x \leq 10$ constraint has persisted into this
new set of constraints.  These constraints are unsatisfiable.

\subsection{Requirements for Constraint Expressions}
\label{sec:constraint-expression-requirements} 

The expressions that define constraints have a number of restrictions.
These will apply to all of the Babelsberg languages.

\begin{enumerate}

\item \label{item:constraint-type} Evaluating the expression that
  defines the constraint should return a boolean.  (This is checked
  dynamically.) 

\item \label{item:side-effects} The constraint expression should be
  free of side
  effects.\footnote{\label{footnote:benign-side-effects}In a practical
    implementation, the programmer might be able to make cautious use of
    benign side effects in a constraint expression, for example, for
    caching or constructing temporary objects that are garbage
    collected before they are visible outside the constraint.  In the
    formal semantics, however, we simply disallow side effects in
    constraint expressions.}

\item \label{item:deterministic} The result of evaluating the block
  should be deterministic.  For example, an expression whose value
  depended on which of two processes happened to complete first
  wouldn't qualify.  (This does not arise in the toy languages here,
  although we do need this restriction for a practical one.)

\end{enumerate}

\subsection{Control Structures}

Babelsberg/Reals includes \verb|if| and \verb|while| control
structures.  These work in the usual way, and allow (for example) a
variable to be incremented only if a test is satisfied, or an
\verb|always| constraint to be conditionally asserted.  The test for
an \verb|if| statement is evaluated, and one or the other branch is
taken --- there is no notion of backtracking to try the other branch.
(Adding Prolog-style backtracking is left for future work --- see
Appendix \ref{sec:adding-solvers}.)  Similarly, a \verb|while|
statement executes the body a fixed number of times --- there is no
possibility of backtracking to execute it a different number of times.

The test in an \verb|if| or \verb|while| statement can use
short-circuit evaluation when evaluating an expression involving
\verb|and| and \verb|or|.  For example, this program  results in
\verb|x| $\mapsto$ 100 (and doesn't get a divide-by-zero error):

\begin{verbatim}
x := 4;
if x=4 or x/0=10
  then x := 100
  else x := 200
\end{verbatim}

For simplicity, our formal rules don't include short-circuit
evaluation --- adding it would be straightforward but would require
additional, not-very-interesting rules.

Constraints with conjunctions or disjunctions are
just turned over to the solver, rather than being evaluated using
short-circuit evaluation.  We could also add \verb|if| expressions to
the language (distinct from \verb|if| statements).  However, since
there is a simple translation from \verb|if| expressions to
conjunctions and disjunctions, we don't include them.  (If we did have
them, they would also need to simply be turned over to the solver.)
For example, the following two constraints are equivalent, and have 
the solution \verb|x| $\mapsto$ 10:

\begin{verbatim}
always if x=4 then x=5 else x=10

always (x=4 and x=5) or (x!=4 and x=10)
\end{verbatim}

In either case, we would turn the entire constraint over to the solver
to find a solution for \verb|x|.

Here is an example of an unsatisfiable constraint of this sort:
\begin{verbatim}
always if x=4 then x=5 else x=4
\end{verbatim}
which is equivalent to:
\begin{verbatim}
always (x=4 and x=5) or (x!=4 and x=4)
\end{verbatim}

\subsection{Adding Other Primitive Types}

It is straightforward to extend Babelsberg/Reals with other primitive
types, such as integers and strings.  When we need to refer to
this language rather than just Babelsberg/Reals we will call it
Babelsberg/PrimitiveTypes.  Note that all the types in
Babelsberg/PrimitiveTypes are atomic --- we don't have recursive types
or types that define values that hold other values (such as records,
arrays, or sets). 

Since we are modeling a dynamically typed language, we need to
consider the case of changing the type of a variable.  Here's a simple
example program that illustrates this.

\begin{verbatim}
x := 5;
x := "Hello";
\end{verbatim}

After the second statement we solve the following constraints:

\begin{tabular}{rl}
weak     & $x = 5$ \\
required & $x = $ \verb|"Hello"|
\end{tabular}

The final result is $\texttt{x} \mapsto \texttt{"Hello"}$.

The changed type might propagate through an \verb|always| constraint:

\begin{verbatim}
x := 5;
y := 10;
always y = x;
x := "Hello";
\end{verbatim}

The final result is that both \verb|x| and \verb|y| are strings.

There is also a potential interaction with overloaded operators.
Suppose that for strings \verb|+| denotes string concatenation:

\begin{verbatim}
x := 5;
y := 10;
always y = x+x;
x := "Hello";
\end{verbatim}

After the \verb|always| statement we have 
$\texttt{x} \mapsto \texttt{5}$, $\texttt{y} \mapsto \texttt{10}$.
Then after the final statement, we have the constraints

\begin{tabular}{rl}
weak     & $x=5$ \\
weak     & $y=10$ \\
required & $y=x+x$ \\
required & $x=$ \texttt{"Hello"}
\end{tabular}

This has the solution $\texttt{x} \mapsto \texttt{"Hello"}$, 
$\texttt{y} \mapsto \texttt{"HelloHello"}$.

Non-determinism can arise for primitive types as well as values:

\begin{verbatim}
x := 3;
always weak x=5;
always weak x="hello";
\end{verbatim}

Here the solver is free to choose either an integer or a string for
\verb|x|.  This is strange, and also not really in keeping with the
spirit of the goals for our languages with respect to constraints and
object identity (Section~\ref{sec:taming-identity-constraints}).
However, it is just an artifact of the semantic rules we are using for
Babelsberg/PrimitiveTypes --- for simplicity, in these rules we don't
distinguish among primitive types, with the consequence that the above
program is legal.  When we get to our final goal (the
Babelsberg/Objects language) the corresponding program will be
disallowed due to the automatic boxing and unboxing of primitive types
in the formal semantic rules.

\subsection{Formalism}
\label{sec:int-formalism}

We present the formal semantics of Babelsberg/PrimitiveTypes. 

\subsubsection{Syntax}
\label{sec:babelsberg-primitivetypes-syntax}


\begin{tt}
\begin{tabular}{lrrl}
{\rm Statement} & s & ::= & skip $\mid$ x := e $\mid$ always C $\mid$ once C $\mid$ s;s \\
  & & & $\mid$ if e then s else s $\mid$ while e do s \\
{\rm Constraint} & C & ::= & $\rho$ e $\mid$ C $\wedge$ C\\
{\rm Expression} & e & ::= & c $\mid$ x $\mid$ e $\oplus$ e  $\mid$ e $\olessthan$ e $\mid$ e $\ovee$ e\\
{\rm Constant} & c & ::= & true $\mid$ false $\mid$ {\rm base type constants}\\
{\rm Variable} & x & ::= & {\rm variable names} \\\\

{\rm Value} & v & ::= & c
\end{tabular}
\end{tt}

The language includes a set of boolean and base type constants (e.g., reals), ranged over by metavariable {\tt c}. A
finite set of operators on expressions is ranged over by $\oplus$. For
booleans this set includes no operations, but it does include
operations on the reals such as {\tt +} and {\tt *}.
The symbol $\olessthan$ ranges over a set of {\em predicate} operators
({\tt $=$} and {\tt $\neq$} for booleans, $\leq$, $<$, $=$, and so on
for reals.)  These test properties of base types in the language. The
symbol $\ovee$ ranges over a set of logical operators for combining
boolean expressions (e.g., {\tt $\wedge$}, {\tt $\vee$}).  The
predicate operators are assumed to include at least an equality
operator {\tt =} for each primitive type in the language, and the
logical operators are assumed to include at least conjunction {\tt $\wedge$}.
The syntax of this language does have some limitations as compared
with that of a practical language --- for example, there are only
binary operators (not unary or ternary), and the result must have the
same type as the arguments.  We make these simplifications since the
purpose of Babelsberg/PrimitiveTypes is to elucidate the semantics of
such languages as a step toward Babelsberg/Objects, rather than to
specify a real language.

For constraints, the symbol $\rho$ ranges over a finite and totally
ordered set of constraint {\em priorities} and is assumed to include a
bottom element {\em weak}
% , an intermediate element {\em strong},
and a top element {\em required}. While syntax requires the priority to be
explicit, for simplicity we sometimes omit it in this semantics. A
constraint with no explicit priority implicitly has the priority {\em
  required}.  Finally, for simplicity we do not model read-only
annotations in the formal semantics.
%% commented out 'strong' since I don't think we make use of it in the
%% formal semantics

The syntax is thus that of a simple, 
standard imperative language except for the
{\tt always} and {\tt once} statements, which declare constraints.
An {\tt always} constraint must
hold for the rest of the programs execution, whereas a {\tt once}
constraint is satisfied by the solver and then retracted. 
%
Note that for simplicity this semantics implicitly gets stuck whenever
the solver cannot satisfy a constraint, either due to an unsatisfiable
constraint or due to the solver being unable to determine whether the
constraint is satisfiable. In a practical implementation, we would
likely want to differentiate between 
these cases, since it's useful if we can
inform the programmer that the constraints are truly not satisfiable.
We could also add standard exception handling to remove the
unsatisfiable or unknown constraint and continue, but omit this here
for simplicity.

\subsubsection{Semantics}


The semantics is defined by several judgments, defined below.  These
judgments depend on the notion of an {\em environment}, which is a
partial function from program variables to program
values. Metavariable {\tt E} ranges over environments. When convenient
we also view an environment as a set of (program variable, program
value) pairs.  For each operator {\tt o} in the language we assume the
existence of a corresponding semantic function denoted
$\llbracket {\tt o} \rrbracket$.

\judgment{\eval{E}{e}{v}}

``Expression {\tt e} evaluates to value {\tt v} in the context of environment {\tt
  E}.''

\inferax{E-Const}{\eval{E}{c}{c}}
\inferrule{E-Var}{
  {\tt E(x) = v}
}{
  \eval{E}{x}{v}
}


% \inferrule{E-ReadOnly}{
%   {\tt E(x) = v}
% }{
%   \eval{E}{x?}{v}
% }

\inferrule{E-Op}{
  \eval{E}{e$_1$}{v$_1$} \andalso
  \eval{E}{e$_2$}{v$_2$} \andalso
  {\tt v$_1$ $\llbracket \oplus \rrbracket$ v$_2$} = {\tt v}
}{
  \eval{E}{e$_1$ $\oplus$ e$_2$}{v}
}

\inferrule{E-Compare}{
  \eval{E}{e$_1$}{v$_1$} \andalso
  \eval{E}{e$_2$}{v$_2$} \andalso
  {\tt v$_1$ $\llbracket \olessthan \rrbracket$ v$_2$} = {\tt v}
}{
  \eval{E}{e$_1$ $\olessthan$ e$_2$}{v}
}

\inferrule{E-Combine}{
  \eval{E}{e$_1$}{v$_1$} \andalso
  \eval{E}{e$_2$}{v$_2$} \andalso
  {\tt v$_1$ $\llbracket \ovee \rrbracket$ v$_2$} = {\tt v}
}{
  \eval{E}{e$_1$ $\ovee$ e$_2$}{v}
}

\inferrule{E-CombineShortCircuit}{
  \eval{E}{e$_1$}{v$_1$} \andalso
  v$_1$ = {\tt true} \andalso
  {\tt v$_1$ $\llbracket \ovee \rrbracket$ v$_2$} = {\tt v}
}{
  \eval{E}{e$_1$ $\ovee$ e$_2$}{v}
}


\judgment{\solve{E}{C}}

This judgment represents a call to the constraint solver, which we
treat as a black box.  The proposition \mbox{\solve{E}{C}}
denotes that environment {\tt E} is a {\em solution} to the constraint
{\tt C} (and further one that is optimal according to the solver's
semantics, as discussed earlier).  

\judgment{\stay{E}{C}}

This judgment defines how to translate an environment into a
source-level ``stay'' constraint.

\inferax{StayEmpty}{\stay{$\emptyset$}{true}}

\inferrule{StayOne}{${\tt E}({\tt x}) = {\tt v}$ \andalso ${\tt E}_0 =
  {\tt E} \setminus \{({\tt x},{\tt v})\}$
  \andalso \stay{E$_0$}{C$_0$} \andalso {\tt C} = {\tt C$_0$}
  $\wedge$ {\tt weak x=v}}
{\stay{E}{C}}

\judgment{\step{E}{C}{s}{E$'$}{C$'$}}

``Execution starting from configuration {\tt <E|C|s>} ends in state {\tt <E$'$|C$'$>}.''

A ``configuration'' defining the state of an execution includes a
concrete context, represented by the environment, a symbolic context,
represented by the constraint, and a statement to be executed.  The
environment and statement are standard, while the constraint is not
part of the state of a computation in most languages.  Intuitively,
the environment comes from constraint solving during the evaluation of
the immediately preceding statement, and the constraint records the
\verb|always| constraints that have been declared so far during
execution. Note that our execution implicitly gets stuck if the solver
cannot produce a model.

\inferrule{S-Asgn}{
  \eval{E}{e}{v} \andalso
  \stay{E}{C$_s$} \andalso
  \solve{E$'$}{(C $\wedge$ C$_s$ $\wedge$ x = v)}
}{
  \step{E}{C}{x := e}{E$'$}{C}
}

\inferrule{S-Once}{
  \stay{E}{C$_s$} \andalso
  \solve{E$'$}{(C $\wedge$ C$_s$ $\wedge$ C$_0$)}
}
{\step{E}{C}{once C$_0$}{E$'$}{C}}

\inferrule{S-Always}{
  \step{E}{C}{once C$_0$}{E$'$}{C} \andalso {\tt C$'$ = C $\wedge$ C$_0$}
}{
  \step{E}{C}{always C$_0$}{E$'$}{C$'$}
}

\inferax{S-Skip}{
  \step{E}{C}{skip}{E}{C}
} 

\inferrule{S-Seq}{
  \step{E}{C}{s$_1$}{E$'$}{C$'$} \andalso
  \step{E$'$}{C$'$}{s$_2$}{E$''$}{C$''$}
}{
  \step{E}{C}{s$_1$;s$_2$}{E$''$}{C$''$}
}

\inferrule{S-IfThen}{
  \eval{E}{e}{true} \andalso
  \step{E}{C}{s$_1$}{E$'$}{C$'$}
}{
  \step{E}{C}{if e then s$_1$ else s$_2$}{E$'$}{C$'$}%
}

\inferrule{S-IfElse}{
  \eval{E}{e}{false} \andalso
  \step{E}{C}{s$_2$}{E$'$}{C$'$}
}{
  \step{E}{C}{if e then s$_1$ else s$_2$}{E$'$}{C$'$}%
}

\inferrule{S-WhileDo}{
  \eval{E}{e}{true} \andalso
  \step{E}{C}{s}{E$'$}{C$'$} \andalso
  \step{E$'$}{C$'$}{while e do s}{E$''$}{C$''$}
}{
  \step{E}{C}{while e do s}{E$''$}{C$''$}
}

\inferrule{S-WhileSkip}{
  \eval{E}{e}{false}
}{
  \step{E}{C}{while e do s}{E}{C}
}


\section{Babelsberg/Records}
\label{sec:babelsberg-records}

For this next language, we augment Babelsberg/PrimitiveTypes with
immutable records.  This language is interesting on its own as an
expository language and as a step toward Babelsberg/Objects; but we
will also continue to use these immutable records in the formal
semantics for Babelsberg/Objects to represent \emph{value classes}
(see Section \ref{sec:babelsberg-objects}).

Records are written as lists of name/value pairs in curly braces:
\begin{verbatim}
  {x:5, y:10}
\end{verbatim}

There is no notion of object identity for Babelsberg/Records --- we can
test whether two records are equal, but whether or not they are
identical would be an implementation issue and not part of the semantics.

The syntax is extended to include record constructors and field
access.  Here are examples of expressions, assignment statements, and
constraints involving records:

\begin{verbatim}
p := {x:2, y:5};   /* assign a record to p */
a := p.x;          /* access a field of a record and assign it to a variable */
q := p;            /* q is now a copy of p (or maybe it's shared; we can't tell
                      the difference) */
always p.x = 100;  /* a constraint on a record field */
always q = p;      /* an equality constraint between two records */
always q.y = 20;
\end{verbatim}

Note the difference between the assignment \verb|q:=p| and the
constraint \verb|always q=p|.  After the \verb|always p.x=100|
constraint, we have $\texttt{p} \mapsto \texttt{\{x:100, y:5\}}$ but
$\texttt{q} \mapsto \texttt{\{x:2, y:5\}}$ --- at this point \verb|p|
and \verb|q| are unrelated, so adding the constraint on \verb|p.x| had no
effect on \verb|q|.  However, after the final constraint, we have
$\texttt{p} \mapsto \texttt{\{x:100, y:20\}}$ and
$\texttt{q} \mapsto \texttt{\{x:100, y:20\}}$, since \verb|p| and
\verb|q| are now constrained to be equal.

The solver must now also handle records.  To tame the power of the
solver so that it does not (for example) invent new fields for
records, we add \emph{structural compatibility checks} on constraints.
These structural compatibility checks are assertions that are checked
dynamically before sending the constraints involving records to the
solver, for example, checking whether a variable is bound to a record,
and whether the record has the necessary fields.  While these
assertions are checked, unlike constraints the system will
never change anything to enforce them --- if one is violated it's just
an error.  Instead, the programmer must ensure that a record with the
expected fields is first assigned to a variable used in record
constraints, just as a programmer would need to ensure that a record
with the expected fields was assigned to a record-valued variable in a
standard language.

Here are a few examples of structural compatibility checks.

\begin{verbatim}
p := {x:2, y:5};
always p.x = 100;
\end{verbatim}
The structural compatibility check is that \verb|p| is a record that
has an \verb|x| field, which succeeds.

The following program is OK --- just as in Babelsberg/PrimitiveTypes,
we can change the type of a variable using an assignment.

\begin{verbatim}
a := {x:1};
a := {y:10};
\end{verbatim}

However, in contrast to Babelsberg/PrimitiveTypes, the following
program fails the structural compatibility checks --- only an
assignment can change the type of a variable.

\begin{verbatim}
a := {x:1};
once a = {y:10};
\end{verbatim}

This program fails the structural compatibility checks as well:
\begin{verbatim}
a := {x:1};
b := {x:1};
always a=b;
a := {x:1, y:10};
\end{verbatim}
Here, assigning a record with a different structure to \verb|a| is OK on its
own, but the \verb|always| constraint would also require \verb|b| to
change.  It's a bit weird that this behavior is different from the
behavior with primitive types (in which a similar program was OK).
However, once we introduce object identity, we will be able to have
changes to the types of variables ripple through the system via
identity constraints (but it will need to be via identity constraints
and not value constraints).

A few more examples:

These two programs fail, because as in Babelsberg/PrimitiveTypes only
assignment can initialize variables.
\begin{verbatim}
once a = {y:10};
\end{verbatim}

\begin{verbatim}
a := {y:10};
always b=a;
\end{verbatim}

The following program also fails the structural compatibility checks, since
the \verb|always| constraint expects \verb|p| to have a \verb|y| field
but it doesn't:
\begin{verbatim}
p := {x:2};
always p.y = 100;
\end{verbatim}

This program fails as well, since constraining a record to be equal to
a number fails the compatibility check:
\begin{verbatim}
p := {x:2};
always p = 5;
\end{verbatim}

The following program passes the structural compatibility checks, but
ends up with an unsatisfiable constraint error, since we are requiring
that \verb|p.x| be both 100 and 2:
\begin{verbatim}  
p := {x:0, y:0};
always p.x = 100;
p := {x:2, y:5};
\end{verbatim}

However, the following program is OK:
\begin{verbatim} 
p := {x:0};
always medium p = {x:3};
always medium p = {x:4}
\end{verbatim}
The solution will be that \verb|p| is a record with a single field
\verb|x| bound to a number, with the exact value depending on the
solver (3, 4, or 3.5 being the most reasonable possibilities).  Note
that the solver may need to adjudicate among solutions that have
different primitive types in a given field, but the structural
compatibility checks ensure that it doesn't need to decide between
two records with e.g.\ different numbers of fields.

Here's another example that fails a structural compatibility check:
\begin{verbatim}
a := {x:0};
b := {y:5};
always a=b
\end{verbatim}

This is the case even though there actually are records that would
satisfy the required \verb|a=b| constraint --- the issue is that the assignment
to \verb|a| leaves it with a record type that has a single \verb|x|
field, the assignment to \verb|b| leaves it with a record type that
has a single \verb|y| field, and the structural compatibility check
prevents the solver from changing either of these types.  This
illustrates one aspect of using the structural compatibility checks to
tame the solver --- one could otherwise imagine the solver coming up
with the solution 
$\texttt{a} \mapsto \texttt{\{x:0, y:5\}}$ and
$\texttt{b} \mapsto \texttt{\{x:0, y:5\}}$, or even
$\texttt{a} \mapsto \texttt{\{x:0, y:5, z:10\}}$ and
$\texttt{b} \mapsto \texttt{\{x:0, y:5, z:10\}}$.  See Appendix 
\ref{sec:structural-compatibility-check-alternatives} for more about
this issue.

\subsection{Formalism}

\subsubsection{Syntax}

The syntax from Section~\ref{sec:babelsberg-primitivetypes} is
augmented now to support records and the ability to access fields of a
record:

\begin{tt}
\begin{tabular}{lrrl}
{\rm Expression} & e & ::= & $\cdots$ $\mid$
\{l$_1$:e$_1$,$\ldots$,l$_n$:e$_n$\} $\mid$ e.l \\
{\rm Label} & l & ::= & {\rm record label names} \\\\
{\rm Value} & v & ::= & c $\mid$ \{l$_1$:v$_1$,$\ldots$,l$_n$:v$_n$\}
\end{tabular}
\end{tt}

We now assume the solver ``understands'' records and record operations
directly.  We assume the equality operator {\tt =} can be used to
compare two records for logical equality, but no other operators apply
to record values.

\subsubsection{Semantics}


\judgment{\eval{E}{e}{v}}

Expression evaluation is updated to support records. 

\inferrule{E-Rec}{
  \eval{E}{e$_1$}{v$_1$} $\cdots$ \eval{E}{e$_n$}{v$_n$}
}{
  \eval{E}{\{l$_1$:e$_1$,$\ldots$,l$_n$:e$_n$\}}{\{l$_1$:v$_1$,$\ldots$,l$_n$:v$_n$\}}
}

\inferrule{E-Field}{
  \eval{E}{e}{\{l$_1$:v$_1$,$\ldots$,l$_n$:v$_n$\}}
  \andalso 1 $\leq$ $i$ $\leq$ $n$
}{
  \eval{E}{e.l$_i$}{v$_i$}
}


\judgment{\tc{E}{e}{T}}

\judgment{\tcC{E}{C}}

``Expression {\tt e} has type {\tt T} in the context of environment
{\tt E}.''

``Constraint {\tt C} is well formed in the context of environment
{\tt E}.''

We use a notion of typechecking to prevent undesirable non-determinism
in constraints.  Specifically, we want constraint solving to preserve
the structure of the values of variables, changing only the underlying
primitive data as part of a solution, in support of the goals listed
in Section \ref{sec:taming-identity-constraints}.  We formalize our
notion of structure through a simple syntax of types:

\begin{tt}
\begin{tabular}{lrrl}
{\rm Type} & T & ::= & PrimitiveType $\mid$ \{l$_1$:T$_1$,$\ldots$,l$_n$:T$_n$\}
\end{tabular}
\end{tt}

The typechecking rules are mostly standard.  However, we check
expressions dynamically just before constraint solving, so we
typecheck in the context of a runtime environment, which is somewhat unusual.

\inferax{T-Constant}{
  \tc{E}{c}{PrimitiveType}
}

\inferrule{T-Variable}{
  \eval{E}{x}{v} \andalso \tc{E}{v}{T}
}{
  \tc{E}{x}{T}
}

% \inferrule{T-ReadOnly}{
%   \tc{E}{x}{T}
% }{
%   \tc{E}{x?}{T}
% }

\inferrule{T-Op}{
  \tc{E}{e$_1$}{PrimitiveType} \andalso \tc{E}{e$_2$}{PrimitiveType}
}{
  \tc{E}{e$_1$ $\oplus$ e$_2$}{PrimitiveType}
}

\inferrule{T-Compare}{
  \tc{E}{e$_1$}{T} \andalso \tc{E}{e$_2$}{T}
}{
  \tc{E}{e$_1$ $\olessthan$ e$_2$}{PrimitiveType}
}

The above rule ensures that two records can only be compared for
equality when they have the same structure.  This is necessary to
ensure that the solver will never have to invent record fields in
order to satisfy a record equality constraint.  

\inferrule{T-Combine}{
  \tc{E}{e$_1$}{PrimitiveType} \andalso \tc{E}{e$_2$}{PrimitiveType}
}{
  \tc{E}{e$_1$ $\ovee$ e$_2$}{PrimitiveType}
}

\inferrule{T-Rec}{
  \tc{E}{e$_1$}{T$_1$} $\cdots$ \tc{E}{e$_n$}{T$_n$}
}{
  \tc{E}{\{l$_1$:e$_1$,$\ldots$,l$_n$:e$_n$\}}{\{l$_1$:T$_1$,$\ldots$,l$_n$:T$_n$\}}
}

\inferrule{T-Field}{
  \tc{E}{e}{\{l$_1$:T$_1$,$\ldots$,l$_n$:T$_n$\}}
  \andalso $1 \leq i \leq n$
}{
  \tc{E}{e.l$_i$}{T$_i$}
}

The above rule ensures that a field of a record can only be referenced
when it already exists in the record. This is necessary to ensure that
the solver will never have to invent record fields in order to satisfy
a constraint containing field accesses.


The following two rules simply ensure that a constraint is well typed.

\inferrule{T-Priority}{\tc{E}{e}{T}}{\tcC{E}{$\rho$ e}}

\inferrule{T-Conjunction}{
  \tcC{E}{C$_1$} \andalso
  \tcC{E}{C$_2$}
}{
  {\tcC{E}{C$_1$ $\wedge$ C$_2$}}
}

%% TODO LATER:
%% \tdm{SHOW: Prove that constraint solving preserves typing, at least
%%   for ``once'' and ``always'' constraints -- not for assignments,
%%   since there we explicitly allow changing types (see below).  And
%%   that assignment statements can only change the type of the variable
%%   that was assigned.  This would be a nice way to validate our choice
%%   of typing rules.}

\judgment{\step{E}{C}{s}{E$'$}{C$'$}}

The semantics for executing statements is essentially identical to what we had
before, except that now we typecheck constraints before we solve
them.  Implicitly we get ``stuck'' if such a check fails.
In a practical language, this could be extended to
generate an exception instead.  We only show the modified rules
below.  

\inferrule{S-Asgn}{
  \eval{E}{e}{v} \andalso \tcC{E[x $\mapsto$ v]}{C} \andalso
  \stay{E}{C$_s$} \andalso
  \solve{E$'$}{(C $\wedge$ C$_s$ $\wedge$ x = v)}
}{
  \step{E}{C}{x := e}{E$'$}{C}
}

The second premise above is necessary so that we don't pass the solver
an ill-typed set of constraints.  Specifically, since the assignment
can change the type of {\tt x}, we have to make sure that the global
constraints in scope, {\tt C}, are well typed with respect to this new
type.  Note that the stay constraint for {\tt x} within {\tt C$_s$} may not be well
typed with respect to the new type of {\tt x}, but that's OK since the
required constraint {\tt x = v} will take precedence.  The syntax {\tt
  E[x $\mapsto$ v]} denotes the environment identical to {\tt E} but
with {\tt x} mapped to the value {\tt v}.

\inferrule{S-Once}{\tcC{E}{C$_0$} \andalso
  \stay{E}{C$_s$} \andalso
  \solve{E$'$}{(C $\wedge$ C$_s$ $\wedge$ C$_0$)}
}
{\step{E}{C}{once C$_0$}{E$'$}{C}}

\inferrule{S-Always}{
\step{E}{C}{once C$_0$}{E$'$}{C} \andalso
{\tt C$'$ = C $\wedge$ C$_0$}
}{
  \step{E}{C}{always C$_0$}{E$'$}{C$'$}
}

\subsection{Adding Mutable Records}

It would be easy to extend Babelsberg/Records
 to allow mutable records.  Syntactically, we would simply
allow field accesses as l-values, e.g.,
\begin{verbatim}
 p := {x:0, y:0};
 p.x := 100;
\end{verbatim}

After the second assignment, we would have 
\verb|p| $\mapsto$ \verb|{x:100, y:0}|.  This doesn't add any
particular complications to the semantics.

However, the only additional feature provided by such an extension
would be the syntax allowing field accesses as l-values --- we can
always convert such a program into one that only has immutable records
by making a new record whose fields are simply copied from those in
the old record, except for the field being assigned.  For example, the
above program is equivalent to the following program in
Babelsberg/Records:
\begin{verbatim}
 p := {x:0, y:0};
 p := {x:100, y:p.y};
\end{verbatim}

In any case there still would be no notion of object identity.  Consider:
\begin{verbatim}
 p := {x:0, y:0};
 q := p;
 p.x := 100;
\end{verbatim}

After executing the three statements \verb|p| $\mapsto$ \verb|{x:100, y:0}|,
but \verb|q| $\mapsto$ \verb|{x:0, y:0}|.

Mutable records would thus simply provide a syntactic convenience,
rather than some new capability.  We would still want immutable
records as well (which we will use to model instances of value classes
in the formal semantics for Babelsberg/Objects), so this would also
add clutter to the languages.  We therefore don't include Babelsberg
with mutable records as a separate language.

\section{Babelsberg/UID}
\label{sec:babelsberg-uid}

Babelsberg/UID adds a number of features to the language. As another step
toward representing objects, we augment Babelsberg/PrimitiveTypes with
records that live on the heap, and that are mutable, have an identity,
and can be aliased.  As with Babelsberg/Records these are represented
as lists of name/value pairs in curly braces:

\begin{verbatim}
  {x:5, y:10}
\end{verbatim}

However, to emphasize that we now allocate records with object
identity on the heap, we use the keyword \verb|new| when creating
one. The syntax includes field accesses as both r-values and l-values.

In all the examples in this section, we will only use records with
object identity.  However, later, when we are formalizing
Babelsberg/Objects, we will use both kinds of records: records with
object identity to model objects that have identity, and immutable
records (as in Babelsberg/Records) to model instances of value classes
(which are immutable objects for which object identity is not
significant).  It should be clear from context when we use the word
``record'' what kind we mean; but when it's necessary to distinguish
them, we'll use the terms ``i-record'' and ``uid-record'' for
immutable records and records with object identity respectively.

Only uid-records live on the heap --- primitive values and i-records do
not.  (It would be possible to store all data on the heap, but we
elected not to, since having everything on the heap makes the
descriptions more complex --- there would always be a level of
indirection to get to data.  If one wants the effect of storing an
integer or boolean or i-record on the heap, it is easy to simulate
this by constructing a uid-record that has a single field that holds
the integer or boolean or i-record.)

A field of a uid-record holds either a primitive type, a i-record,
or a reference to another uid-record.  There is no syntax for creating
a nested uid-record directly --- nested records have to be constructed
with references, so each record \emph{must} have a variable that
refers to it, even if it is only used in a nested structure. This
simplifies the semantics.  (Nested uid-records could be supported
directly with just a source code transformation.  In any case, when we
get to Babelsberg/Objects, we will remove this restriction and allow
expressions that do involve creating new objects.)

Since object identity is now significant, we add identity constraints
to the language (following Smalltalk syntax, 
written \verb|==|, in contrast to \verb|=| for
equality constraints).  For records \verb|p| and \verb|q|, if \verb|p|
and \verb|q| are identical, they must also be equal, but the converse
is not necessarily true --- if \verb|p| and \verb|q| are equal, they
might or might not be identical.

Unlike Babelsberg/PrimitiveTypes and Babelsberg/Records, the weak
stays on variables are identity constraints rather than equality
constraints.  For uid-records, such a variable continues to refer to
the same object unless reassigned.  This is a direct consequence of
the weak stays referring to the references the variables hold, not the
records on the heap.  In addition, there are also weak stay
constraints on the values in the record fields.

For simplicity, we also allow identity tests and constraints on
primitive types and i-records.  (Otherwise we would need two different
translations for assignment.)  Again for simplicity, two instances of
a primitive type or i-record are identical iff they are equal.
(Allowing them to be equal but not identical would seem to imply
storing them on the heap, which would complicate the semantics.)

Here are some examples involving uid-record fields and constraints:
\begin{verbatim}
p := new {x:2, y:5}; /* create a record on the heap and save the reference in p */
a := p.x;            /* read a field of a record and assign it to a variable */
p.x := 6;            /* assign 6 to the field p.x */
always p.x = 100;
\end{verbatim}

As with Babelsberg/Records, there is no static type checking ---
checking is all done dynamically.  As before, our model includes
structural compatibility checks on constraints, which tames the solver so
that it doesn't do such things as inventing new records or adding fields
to a record.  (The only way to create a new record is with an explicit
\verb|new| expression in the program, and this also defines all of the
fields that it has and will ever have.)

After executing the right hand side of the first statement, the heap
is updated with a new object and a reference $r$ to point to it. The
assignment and stay constraints then desugar into:

\begin{tabular}{rl}
required & $p = r$ \\
required & $H(r) = \{x:x_r, y:y_r\}$\\
weak & $x_r = 2$\\
weak & $y_r = 5$\\
\end{tabular}

The reference $r$ is a constant. We assume the solver understands
uninterpreted function symbols.  The uninterpreted function symbol $H$
in the second constraint above is used to represent the heap's mapping
from references to their contents; the constraint forces the solver to
keep the reference $r$ pointing to a record that has \emph{only} $x$
and $y$ fields. We introduce variables for the values of these fields
and add weak stays so these keep the values with which they were
created -- these can always be satisfied, because the variables are
fresh.

After the second statement we have the stay constraints on \verb|p|
and its fields, and the required constraint resulting from the
assignment:

\begin{tabular}{rl}
weak & $p = r$ \\
required & $H(r) = \{x:x_r, y:y_r\}$\\
weak & $x_r = 2$\\
weak & $y_r = 5$\\
required & $a = 2$
\end{tabular}

After the third statement:

\begin{tabular}{rl}
weak & $p = r$ \\
required & $H(r) = \{x:x_r, y:y_r\}$\\
weak & $x_r = 2$\\
weak & $y_r = 5$\\
weak & $a = 2$ \\
required & $p.x = 6$
\end{tabular}


And finally:

\begin{tabular}{rl}
weak & $p = r$ \\
required & $H(r) = \{x:x_r, y:y_r\}$\\
weak & $x_r = 6$\\
weak & $y_r = 5$\\
weak & $a = 2$ \\
required & $p.x = 100$
\end{tabular}

After solving the final set of constraints we have the environment 
\verb|p| $\mapsto$ \verb|r| $\wedge$
\verb|a| $\mapsto$ \verb|2|, as well as the heap \linebreak \verb|r| $\mapsto$
\verb|{x:100, y:5}|.

Here is an example that fails the structural compatibility checks.

\begin{verbatim}
p := new {x:2, y:5};
always p.z = 5;
\end{verbatim}

As before, the system is not allowed to add fields to records, so \verb|p|
would be required to have a \verb|z| field already.

The following example demonstrates one form of aliasing:
\begin{verbatim}
p := new {x:2, y:5};
q := p;
p.x := 100;
q := new {z:10};
p.x := 200;
\end{verbatim}

After the assignment \verb|q:=p|, \verb|q| holds the same reference as
\verb|p|, so the subsequent assignment to \verb|p.x| changes both.
But then we break the alias with another assignment to \verb|q|, so
the second assignment to \verb|p.x| doesn't affect \verb|q| (which by
that point no longer has an \verb|x| field).

There can also be explicit identity constraints:
\begin{verbatim}
p := new {x:2, y:5};
q := p;
always q==p;
q := new {z:10};
\end{verbatim}
After the last statement both \verb|p| and \verb|q| point to the same
\verb|{z:10}| record.  Note that we don't have structural
compatibility checks for re-satisfying identity constraints, so that
the assignment to \verb|q| that changes its structure ripples through
the identity constraint
to change \verb|p| as well.  However, if the \verb|always| constraint
had been a record equality constraint rather than an identity
constraint, the last assignment to \verb|q| would have failed the
structural compatibility check, which would have expected it to be a
record with \verb|x| and \verb|y| fields.

In Section \ref{sec:taming-identity-constraints} we listed a number of
goals for our design to attempt to tame the power of identity
constraints and the solver.  In support of this, we require that
variables be created in an assignment statement --- the above program
would have been illegal if we didn't have the \verb|q:=p| statement
but just tried to create \verb|q| with the \verb|always| constraint.
We also require that a new identity constraint be satisfied at the
time it is created --- otherwise it's an error (modeled in the formal
semantics by the rules getting stuck).  This restriction was obeyed in
the above program.  Here is an example that violates it, and is hence
not legal:
\begin{verbatim}
p := new {x:2};
q := new {y:5};
always q==p;
\end{verbatim}

Without the restriction, after the last statement \verb|p| and
\verb|q| would definitely point to the same record, but it might be
either \verb|{x:2}| or \verb|{y:5}|.  (In any case, \verb|p| and
\verb|q| would still one of the records that was already created by
the \verb|new| expressions --- the solver wouldn't invent a new one ---
but we want to tame the language further.)  

With the restriction, any changes to the types of variables must flow
from an explicit assignment statement --- the effects can still ripple
outwards via identity constraints --- but they are deterministic.  For
the same reason, we disallow disjunctions of identity constraints
(since otherwise there could be programs with multiple correct
solutions to the identity constraints).  It is only updates to values
that can be non-deterministic as the result of constraints.  One item
for future work will be to formalize and prove this statement, which
seems like a useful property for the language.

% The programmer can further restrict what happens with an explicit
% identity constraint by adding a read-only annotation.  For example,
% this program is legal:
% \begin{verbatim}
% p := new {x:2};
% q := p;
% always q==p?;
% p := new {y:10};
% \end{verbatim}

% But this one fails, due to the read-only annotation on \verb|p| in the
% identity constraint:

% \begin{verbatim}
% p := new {x:2};
% q := p;
% always q==p?;
% q := new {y:10};
% \end{verbatim}

% \ab{Do we even want to allow read-only annotations on variables in
%   identity constraints?  They seem potentially useful in a practical
%   language, as a machine-checkable restriction on how the identity
%   constraint is used.  But they aren't essential for the formal
%   semantics.}

As another aspect of taming identity constraints, explicit priorities
on identity constraints are not allowed --- there are only the
implicit weak identity stay constraints on all variables.  (We haven't
found any compelling use cases for programmer-specified soft identity
constraints, and omitting them simplifies reasoning about identity
constraints for the programmer.)  Similarly, we disallow disjunctions
of identity constraints, which would also introduce nondeterminism;
and for simplicity, there are no explicit conjunctions of identity
constraints, only the implicit conjunction resulting from writing
several of them.

This program is thus illegal:

\begin{verbatim}
p := new {x:0};
q := new {x:5};
always medium p.x = 0;
always medium q.x = 5;
always weak p==q
\end{verbatim}

Without the restriction, the programmer would need to interleave
determining values and identities in reasoning about the program's
behavior.  For example, for the above program, to decide that we
should leave the weak identity constraint unsatisfied, we'd need to
first solve the value constraints and then decide that we can't
satisfy the weak identity constraint.

The solver can't spontaneously create new uid-records if they are needed
--- new uid-records can only be created with an explicit assignment with a
\verb|new| on the right hand side.  Consider:

\begin{verbatim}
a := new {x:1};
b := a;
always a.x=1;
always b.x=2;
\end{verbatim}

After \verb|b:=a|, {\tt a} and {\tt b} refer to the same record, but
after the second {\tt always} constraint, the solver would need to
spontaneously create a new record and point {\tt b} at it to satisfy
the constraints.  So this program halts with an unsatisfiable
constraint error.

It also can't switch identities around to satisfy value constraints.
Here is the same program, except that there happens to be a record
with the required value for its \verb|x| field already lying around.
But to no avail: this program halts with an unsatisfiable constraint
error as well, rather than silently
changing \verb|b| to refer to the \verb|{x:2}| record.

\begin{verbatim}
a := new {x:1};
b := a;
c := new {x:2};
always a.x=1;
always b.x=2;
\end{verbatim}

\subsection{Formalism}

Since this version is significantly different than the prior formalisms, 
we present it in its entirety rather than as a delta from those ones.

\subsubsection{Syntax}

The syntax is augmented to create records on the heap with the {\tt
  new} keyword. We also add references as values and an identity
operation {\tt ==} to compare values by identity rather than by
structure (for primitive values, this is the same as {\tt =}). The
arguments to identity operations are not arbitrary expressions but
rather {\em accessors}, which are essentially access paths through a
reference.  This is no loss of expressiveness but is useful for the
translation to the constraint solver.

\begin{tt}
\begin{tabular}{lrrl}
{\rm Statement} & s & ::= & skip $\mid$ a := e $\mid$ a := new o $\mid$ always C $\mid$ once C $\mid$ s;s \\
  & & & $\mid$ if e then s else s $\mid$ while e do s \\
{\rm Constraint} & C & ::= & $\rho$ e $\mid$ C $\wedge$ C\\
{\rm Expression} & e & ::= & c $\mid$ r $\mid$ a $\mid$ e $\oplus$ e  $\mid$ e
$\olessthan$ e $\mid$ e $\ovee$ e $\mid$ I $\mid$ D \\
{\rm Identity} & I & ::= & a == a\\
{\rm Accessor} & a & ::= & x $\mid$ a.l \\
{\rm Record} & o & ::= & \{l$_1$:e$_1$,$\ldots$,l$_n$:e$_n$\} \\
{\rm Dereference} & D & ::= & H(r)\\
{\rm Reference} & r & ::= & {\rm references to heap records} \\\\
{\rm Constant} & c & ::= & true $\mid$ false $\mid$ {\rm base type constants}\\
{\rm Variable} & x & ::= & {\rm variable names} \\
{\rm Label} & l & ::= & {\rm label names} \\\\

{\rm Value} & v & ::= & c $\mid$ r
\end{tabular}
\end{tt}


In the syntax, we treat {\tt H} as a keyword.  Source programs will
not use expressions of the form {\tt H(e)}, but they are introduced as
part of constraints given to the solver, which we assume will treat
{\tt H} as an uninterpreted function, as described earlier.

\subsubsection{Semantics}

In addition to an environment, the semantics now requires a {\em
  heap}, which is a partial function from mutable references to record
``objects'' of the form {\tt \{l$_1$:v$_1$,$\ldots$,l$_n$:v$_n$\}}.
Heaps are ranged over by metavariable {\tt H}\footnote{Not to be
  confused with the keyword {\tt H} in the syntax; it is always clear from
  context which is being used.}. Secondly, in addition
to the normal active constraints, we also now have a number of
required identity constraints that are active. These are represented
by {\tt I}. 

\judgment{\evalH{E}{H}{e}{v}}

``Expression {\tt e} evaluates to value {\tt v} in the context of environment {\tt
  E} and heap {\tt H}.''

Most of the updates to the expression rules just use the new judgment
that includes the heap.  We do not include rules for evaluating
expressions of the form {\tt H(r)}, since they are not intended to
appear in source programs.

\inferax{E-Const}{\evalH{E}{H}{c}{c}}

\inferrule{E-Var}{{\tt E}({\tt x}) = {\tt v}}{\evalH{E}{H}{x}{v}}

\inferrule{E-Field}{
  \evalH{E}{H}{a}{r} \andalso
  {\tt H}({\tt r}) = {\tt \{l$_1$:v$_1$,$\ldots$,l$_n$:v$_n$\}} \andalso
  1 $\leq$ $i$ $\leq$ $n$
}{
  \evalH{E}{H}{a.l$_i$}{v$_i$}
}


\inferax{E-Ref}{\evalH{E}{H}{r}{r}}

\inferrule{E-Op}{
  \evalH{E}{H}{e$_1$}{v$_1$} \andalso
  \evalH{E}{H}{e$_2$}{v$_2$} \andalso
  {\tt v$_1$ $\llbracket \oplus \rrbracket$ v$_2$} = {\tt v}
}{
  \evalH{E}{H}{e$_1$ $\oplus$ e$_2$}{v}
}

\inferrule{E-Compare}{
  \evalH{E}{H}{e$_1$}{v$_1$} \andalso
  \evalH{E}{H}{e$_2$}{v$_2$} \andalso
  {\tt v$_1$ $\llbracket \olessthan \rrbracket$ v$_2$} = {\tt v}
}{
  \evalH{E}{H}{e$_1$ $\olessthan$ e$_2$}{v}
}

\inferrule{E-Combine}{
  \evalH{E}{H}{e$_1$}{v$_1$} \andalso
  \evalH{E}{H}{e$_2$}{v$_2$} \andalso
  {\tt v$_1$ $\llbracket \ovee \rrbracket$ v$_2$} = {\tt v}
}{
  \evalH{E}{H}{e$_1$ $\ovee$ e$_2$}{v}
}

\infrule[E-IdentityTrue]{
  \mbox{%
    \evalH{E}{H}{a$_1$}{v} \andalso
    \evalH{E}{H}{a$_2$}{v}}
}{\mbox{%
  \evalH{E}{H}{a$_1$ == a$_2$}{true}
}}

\infrule[E-IdentityFalse]{
  \mbox{%
    \evalH{E}{H}{a$_1$}{v$_1$} \andalso
    \evalH{E}{H}{a$_2$}{v$_2$} \andalso {\tt v$_1$} $\neq$ {\tt v$_2$}}
}{\mbox{%
  \evalH{E}{H}{a$_1$ == a$_2$}{false}
}}

\judgment{\tcH{E}{H}{e}{T}}
\judgment{\tcHC{E}{H}{C}}

The typing rules also are augmented with the heap. They are otherwise
unchanged. The additional thing to type are references. We do not add
type rules for identities. This ensures that constraints involving
them do not typecheck, so identity checks cannot occur in ordinary
constraints.  As with the evaluation rules, expressions of the form
{\tt H(r)} do not typecheck.

\inferax{T-Constant}{
  \tcH{E}{H}{c}{PrimitiveType}
}

\inferrule{T-Ref}{
  \tcH{E}{H}{o}{\{l$_1$:T$_1$,$\ldots$,l$_n$:T$_n$\}} \andalso
  {\tt H}({\tt r})={\tt o}
}{
  \tcH{E}{H}{r}{\{l$_1$:T$_1$,$\ldots$,l$_n$:T$_n$\}}
}

\inferrule{T-Variable}{
  \eval{E}{x}{v} \andalso \tcH{E}{H}{v}{T}
}{
  \tcH{E}{H}{x}{T}
}

\inferrule{T-Field}{
  \tcH{E}{H}{a}{\{l$_1$:T$_1$,$\ldots$,l$_n$:T$_n$\}}
  \andalso $1 \leq i \leq n$
}{
  \tcH{E}{H}{a.l$_i$}{T$_i$}
}

% \inferrule{T-ReadOnly}{
%   \tcH{E}{H}{x}{T}
% }{
%   \tcH{E}{H}{x?}{T}
% }

\inferrule{T-Op}{
  \tcH{E}{H}{e$_1$}{PrimitiveType} \andalso \tcH{E}{H}{e$_2$}{PrimitiveType}
}{
  \tcH{E}{H}{e$_1$ $\oplus$ e$_2$}{PrimitiveType}
}

\inferrule{T-Compare}{
  \tcH{E}{H}{e$_1$}{T} \andalso \tcH{E}{H}{e$_2$}{T}
}{
  \tcH{E}{H}{e$_1$ $\olessthan$ e$_2$}{PrimitiveType}
}


\inferrule{T-Combine}{
  \tcH{E}{H}{e$_1$}{PrimitiveType} \andalso \tcH{E}{H}{e$_2$}{PrimitiveType}
}{
  \tcH{E}{H}{e$_1$ $\ovee$ e$_2$}{PrimitiveType}
}



\inferrule{T-Priority}{\tcH{E}{H}{e}{T}}{\tcHC{E}{H}{$\rho$ e}}

\inferrule{T-Conjunction}{
  \tcHC{E}{H}{C$_1$} \andalso
  \tcHC{E}{H}{C$_2$}
}{
  {\tcHC{E}{H}{C$_1$ $\wedge$ C$_2$}}
}

\judgment{\solveH{E}{H}{C}}

This still represents the call to the solver. The solver now
determines values on the heap and in the environment.

%% TODO LATER:
%% \tf{SHOW: The solver won't invent references or change object
%%   structures on the heap. Only fields and variables pointing to
%%   references are allowed to change to point to other references.}

\judgment{\stay{E}{C}}
\judgment{\stay{H}{C}}

These judgments define how to translate the environment and heap into a
source-level ``stay'' constraint. The rules are unchanged for
primitive values that are stored in the environment, but are amended
for the heap. We assume the solver knows about maps to do lookup on
the heap.

\inferax{StayEmpty}{\stay{$\emptyset$}{true}}

\inferrule{StayOne}{${\tt E}({\tt x}) = {\tt v}$ \andalso ${\tt E}_0 =
  {\tt E} \setminus \{({\tt x},{\tt v})\}$
  \andalso \stay{E$_0$}{C$_0$} \andalso {\tt C} = {\tt C$_0$}
  $\wedge$ {\tt weak x=v}}
{\stay{E}{C}}

\inferrule{StayHeap}{
  ${\tt o} = $\{{\tt l$_0$:v$_0$,$\ldots$,l$_n$:v$_n$}\} \andalso
  ${\tt H}({\tt r}) = {\tt o}$ \andalso  
  ${\tt H}_0 = {\tt H} \setminus \{({\tt r},{\tt o})\}$ \andalso
  \stay{H$_0$}{C$_0$}
}{
  \stay{H}{C$_0$ $\wedge$
    H(r) = \{l$_0$:x$_{r0}$,$\ldots$,l$_n$:x$_{rn}$\} $\wedge$
    weak x$_{r0}$ = v$_0$ $\wedge$
    $\ldots$ $\wedge$
    weak x$_{rn}$ = v$_n$ }
}

Note that the constraints involving {\tt H(r)} are {\em required},
which ensures that a reference keeps pointing to the same object.

\judgment{\stepH{E}{H}{C}{I}{s}{E$'$}{H$'$}{C$'$}{I$'$}}

The semantics for executing statements is extended from what we had
before for Babelsberg/Records. 
When updating a variable, we now solve first for the identity
constraints and then for the value constraints. This propagates any
type changes through the identity constraints. We then type the value
constraints against the environment returned by solving for the
identities.

\infrule[S-Asgn]{
  \mbox{
    \evalH{E}{H}{e}{v}}\\
  \mbox{
    \stay{E}{C$_{E_s}$} \andalso
    \stay{H}{C$_{H_s}$} \andalso
    \solveH{E$'$}{H$'$}{(I $\wedge$ C$_{E_s}$ $\wedge$ C$_{H_s}$ $\wedge$ a=v)}}\\
  \mbox{
    \stay{E$'$}{C$_{{E'}_s}$} \andalso
    \stay{H$'$}{C$_{{H'}_s}$} \andalso
    \tcHC{E$'$}{H$'$}{C} \andalso
    \solveH{E$''$}{H$''$}{(C $\wedge$ C$_{{E'}_s}$ $\wedge$ C$_{{H'}_s}$ $\wedge$ a=v)}}
}{\mbox{
  \stepH{E}{H}{C}{I}{a := e}{E$''$}{H$''$}{C}{I}
}}

Implicitly this rule gets stuck if either a) the identity constraints
cannot be solved, b) the value constraints do not type if the
identities are solved, or c) the value constraints cannot be solved. A
practical implementation would add explicit exceptions for these cases
that the programmer could handle.


\infrule[S-AsgnNew]{
  \mbox{
    \evalH{E}{H}{e$_1$}{v$_n$} $\cdots$ \evalH{E}{H}{e$_n$}{v$_n$}} \\
\mbox{
    {\tt H}({\tt r})$\uparrow$ \andalso {\tt H$_0$} = {\tt H}
    $\bigcup$ \{({\tt r}, {\tt
      \{l$_1$:v$_1$,$\ldots$,l$_n$:v$_n$\}})\}} \\
  \mbox{
    \stay{E}{C$_{E_s}$} \andalso
    \stay{H$_0$}{C$_{H_s}$} \andalso
    \solveH{E$'$}{H$'$}{(I $\wedge$ C$_{E_s}$ $\wedge$ C$_{H_s}$ $\wedge$ a=r)}}\\
  \mbox{
    \stay{E$'$}{C$_{{E'}_s}$} \andalso
    \stay{H$'$}{C$_{{H'}_s}$} \andalso
    \tcHC{E$'$}{H$'$}{C} \andalso
    \solveH{E$''$}{H$''$}{(C $\wedge$ C$_{{E'}_s}$ $\wedge$ C$_{{H'}_s}$ $\wedge$ a=r)}}
}{\mbox{
  \stepH{E}{H}{C}{I}{a := new \{l$_1$:e$_1$,$\ldots$,l$_n$:e$_n$\}}{E$''$}{H$''$}{C}{I}
}}

The next rules are essentially unchanged from their previous versions
and just augmented with the heap references.

\inferrule{S-Once}{
  \tcHC{E}{H}{C$_0$} \andalso
  \stay{E}{C$_{E_s}$} \andalso
  \stay{H}{C$_{H_s}$} \andalso
  \solveH{E$'$}{H$'$}{$($C $\wedge$ C$_{E_s}$ $\wedge$ C$_{H_s}$ $\wedge$ C$_0$$)$}
}
{\stepH{E}{H}{C}{I}{once C$_0$}{E$'$}{H$'$}{C}{I}}

\inferrule{S-Always}{
\stepH{E}{H}{C}{I}{once C$_0$}{E$'$}{H$'$}{C}{I}
\andalso {\tt C$'$ = C $\wedge$ C$_0$}
}{
  \stepH{E}{H}{C}{I}{always C$_0$}{E$'$}{H$'$}{C$'$}{I}
}

\inferax{S-Skip}{
  \stepH{E}{H}{C}{I}{skip}{E}{H}{C}{I}
}

\inferrule{S-Seq}{
  \stepH{E}{H}{C}{I}{s$_1$}{E$'$}{H$'$}{C$'$}{I$'$} \andalso
  \stepH{E$'$}{H$'$}{C$'$}{I$'$}{s$_2$}{E$''$}{H$''$}{C$''$}{I$''$}
}{
  \stepH{E}{H}{C}{I}{s$_1$;s$_2$}{E$''$}{H$''$}{C$''$}{I$''$}
}

\inferrule{S-IfThen}{
  \evalH{E}{H}{e}{true} \andalso
  \stepH{E}{H}{C}{I}{s$_1$}{E$'$}{H$'$}{C$'$}{I$'$}
}{
  \stepH{E}{H}{C}{I}{if e then s$_1$ else s$_2$}{E$'$}{H$'$}{C$'$}{I$'$}%
}

\inferrule{S-IfElse}{
  \evalH{E}{H}{e}{false} \andalso
  \stepH{E}{H}{C}{I}{s$_2$}{E$'$}{H$'$}{C$'$}{I$'$}
}{
  \stepH{E}{H}{C}{I}{if e then s$_1$ else s$_2$}{E$'$}{H$'$}{C$'$}{I$'$}%
}

\inferrule{S-WhileDo}{
  \evalH{E}{H}{e}{true} \andalso
  \stepH{E}{C}{H}{I}{s}{E$'$}{H$'$}{C$'$}{I$'$} \andalso
  \stepH{E$'$}{H$'$}{C$'$}{I$'$}{while e do s}{E$''$}{H$''$}{C$''$}{I$''$}
}{
  \stepH{E}{H}{C}{I}{while e do s}{E$''$}{H$''$}{C$''$}{I$''$}
}

\inferrule{S-WhileSkip}{
  \evalH{E}{H}{e}{false}
}{
  \stepH{E}{H}{C}{I}{while e do s}{E}{H}{C}{I}
}

The next rules are additions to create identity constraints. Note that
the above rules for always and once match only constraints that do not
include identity comparisons, because those would not type.  We
require that an identity constraint already be satisfied when it is
asserted, as discussed earlier.  So the following rules simply need to
check the constraint.

\infrule[S-OnceIdentity]{
  \mbox{
    \evalH{E}{H}{a$_0$}{v} \andalso
    \evalH{E}{H}{a$_1$}{v}}
}{\mbox{
    \stepH{E}{H}{C}{I}{once a$_0$ == a$_1$}{E}{H}{C}{I}
}}

\infrule[S-AlwaysIdentity]{
\mbox{
    \stepH{E}{H}{C}{I}{once a$_0$ == a$_1$}{E}{H}{C}{I}
\andalso {\tt I$'$ = I $\wedge$ (a$_0$ == a$_1$)}}
}{\mbox{
    \stepH{E}{H}{C}{I}{always a$_0$ == a$_1$}{E}{H}{C}{I$'$}
}}


% \section{Babelsberg/Accessors}
% \label{sec:babelsberg-accessorss}

% In this language, we add classes, inheritance, methods, and messages.
% This language is an intermediate step towards full
% Babelsberg/Objects. This language adds method invocation and inlining
% of method calls into constraints, but still does not implement full
% constraint construction mode of actual object constraint languages
% including Babelsberg/R, Babelsberg/JS, and Babelsberg/Squeak. Instead,
% this language only allows accessor methods to be used in constraints.

% Since concerns like class and method definitions are orthogonal to
% this semantics, we omit them entirely.

% \subsection{Formalism}

% \subsubsection{Syntax}

% The syntax is augmented to support method invocation. We introduce
% syntax elements for the method body and the formal parameters and {\tt
%   self}, which can be referenced from a method but not assigned
% to. Additionally, we also allow creating new objects as expressions
% now. We remove records values. Records are only used to create a new
% heap structure (we assume class constructors would use this primitive
% to instantiate objects.) The entire syntax is given below.

% \begin{tt}
% \begin{tabular}{lrrl}
% {\rm Statement} & s & ::= & skip $\mid$ L := e $\mid$ always C $\mid$ once C $\mid$ s;s \\
%   & & & $\mid$ if e then s else s $\mid$ while e do s \\
% {\rm Constraint} & C & ::= & e $\mid$ $\rho$ e \\
% {\rm Expression} & e & ::= & c $\mid$ a $\mid$ a? $\mid$ e $\oplus$ e  $\mid$ e $\olessthan$ e $\mid$ e $\ovee$ e $\mid$ I \\
% {\rm Identity} & I & ::= & a == a\\
% {\rm Accessor} & a & ::= & L $\mid$ f $\mid$ a.l(e$_1$,$\ldots$,e$_n$) $\mid$ new \{l$_1$:e$_1$,$\ldots$,l$_n$:e$_n$\} \\
% {\rm L-Value} & L & ::= & x $\mid$ a.l \\
% {\rm Formal} & f & ::= & self $\mid$ {\rm variable names} \\
% {\rm Constant} & c & ::= & true $\mid$ false $\mid$ nil $\mid$ {\rm base type constants}\\
% {\rm Variable} & x & ::= & {\rm variable names} \\
% {\rm Label} & l & ::= & {\rm record label names} \\
% {\rm Reference} & r & ::= & {\rm references to heap records} \\\\
% {\rm Method Body} & b & ::= & s; return e\\

% {\rm Value} & v & ::= & c $\mid$ r\\
% \end{tabular}
% \end{tt}

% Method bodies are assumed to include one or more statements and then
% return an expression. A method body that includes only a return can be
% expressed as {\tt skip; return e}.

% \subsubsection{Semantics}

% \judgment{\evalBig{E}{H}{e}{H}{v}}

% Expression evaluation is changed to account for method invocation and
% object instantiation as expressions. These operations may have an
% effect on the heap, the active constraints, and the identity
% constraints, so this judgment now returns updates to these besides the
% value.

% Evaluation also needs two helpers, one for method lookup and one for
% setting up method scope and evaluating the argument. The precise
% semantics of method lookup are not important, so we just use it as an
% opaque (partial) function from values and labels to statements.

% \judgment{\lookup{v}{l}{b}}

% ``Lookup of method {\tt l} in the object or value {\tt v} returns the
% formal parameter names {\tt f$_1$} through {\tt f$_n$} and the method
% body {\tt b}''

% \judgment{\enter{E}{H}{C}{I}{v}{E$_m$}{H$'$}{C$'$}{I$'$}}

% ``Invoking a method on {\tt v} with argument names {\tt f$_1$} through
% {\tt f$_n$} and arguments {\tt e$_1$} through {\tt e$_n$} constructs
% the method scope {\tt E$_m$} and may have side effects.''

% \infrule[Enter]{
%   \mbox{
%   \evalBigC{E}{H}{C}{I}{e$_1$}{H$_1$}{C$_1$}{I$_1$}{v$_1$}
%   $\cdots$
%   \evalBigC{E}{H$_{n-1}$}{C$_{n-1}$}{I$_{n-1}$}{e$_n$}{H$_n$}{C$_n$}{I$_n$}{v$_n$}}\\
% \mbox{
%   {\tt E$_m$} = \{({\tt self}, {\tt v}),({\tt f$_1$}, {\tt v$_1$}),$\cdots$,({\tt f$_n$}, {\tt v$_n$})\}}
% }{\mbox{
%   \enter{E}{H}{C}{I}{v}{E$_m$}{H$_n$}{C$_n$}{I$_n$}
% }}

% \infrule[E-Call]{\mbox{
%   \evalBigC{E}{H}{C}{I}{a}{H$_0$}{C$_0$}{I$_0$}{v} \andalso
%   \lookup{v}{l}{s; return e}}\\
% \mbox{
%   \enter{E}{H$_0$}{C$_0$}{I$_0$}{v}{E$_m$}{H$_n$}{C$_n$}{I$_n$}}\\
% \mbox{\stepH{E$_m$}{H$_n$}{C$_n$}{I$_n$}{s}{E$'$}{H$'$}{C$'$}{I$'$} \andalso
%   \evalBigC{E$'$}{H$'$}{C$'$}{I$'$}{e}{H$''$}{C$''$}{I$''$}{v$_r$}}
% }{\mbox{
%   \evalBigC{E}{H}{C}{I}{a.l(e$_1$,$\ldots$,e$_n$)}{H$''$}{C$''$}{I$''$}{v$_r$}
% }}

% \infrule[E-New]{\mbox{
%   \evalBigC{E}{H}{C}{I}{e$_1$}{H$_1$}{C$_1$}{I$_1$}{v$_1$} $\cdots$ 
%   \evalBigC{E$_{n-1}$}{H$_{n-1}$}{C$_{n-1}$}{I$_{n-1}$}{e$_n$}{H$_n$}{C$_n$}{I$_n$}{v$_n$}}\\
% \mbox{
%   {\tt H$'$} = ({\tt H$_n$} $\bigcup$ \{({\tt r}, {\tt \{l$_1$:v$_1$,$\ldots$,l$_n$:v$_n$\}})\})}
% }{\mbox{
%   \evalBigC{E}{H}{C}{I}{new \{l$_1$:e$_1$,$\ldots$,l$_n$:e$_n$\}}{H$'$}{C$_n$}{I$_n$}{r}
% }}


% \judgment{\tcH{E}{H}{e}{T}}
% \judgment{\tcH{E}{H}{s}{T}}
% \judgment{\tcHC{E}{H}{C}}

% The typing rules are unchanged. In particular, even though we allow
% method calls in expressions and thus in constraints syntactically,
% these do not type. Instead, our rules for creating constraints given
% below inline method invocations.

% \judgment{\solveH{E}{H}{C}}

% For simplicity of this semantics we assumed that the solver natively
% supports records and maps. We do not assume, however, that the solver
% understands methods, which can now be part of constraint
% expressions. This requires us to essentially inline methods before
% passing constraints to the solver.

% \judgment{\stay{E}{C}}
% \judgment{\stay{H}{C}}

% The rules for creating stay constraints are unchanged.

% \judgment{\inline{E}{H}{e}{H$'$}{e$'$}}
% \judgment{\inlineS{E}{H}{s}{E$'$}{H$'$}{s$'$}}

% ``Inlining expression {\tt e} may update the heap and otherwise turns
% into expression {\tt e$'$}''

% ``Inlining statement {\tt s} may update the environment and heap and
% otherwise turns into expression {\tt s$'$}''

% \inferax{I-Const}{\inline{E}{H}{c}{H}{c}}

% \inferax{I-Var}{\inline{E}{H}{x}{H}{x}}

% \inferrule{I-Field}{
%   \inline{E}{H}{a}{H$'$}{a$'$}
% }{
%   \inline{E}{H}{a.l}{H}{a$'$.l}
% }

% \inferax{I-ReadOnly}{\inline{E}{H}{a?}{H}{a?}}

% \inferax{I-Ref}{\inline{E}{H}{r}{H}{r}}

% \inferrule{I-Op}{
%   \inline{E}{H}{e$_1$}{H$'$}{e$_a$} \andalso
%   \inline{E}{H$'$}{e$_2$}{H$''$}{e$_b$}
% }{
%   \inline{E}{H}{e$_1$ $\oplus$ e$_2$}{H$''$}{e$_a$ $\oplus$ e$_b$}
% }

% \inferrule{I-Compare}{
%   \inline{E}{H}{e$_1$}{H$'$}{e$_a$} \andalso
%   \inline{E}{H$'$}{e$_2$}{H$''$}{e$_b$}
% }{
%   \inline{E}{H}{e$_1$ $\olessthan$ e$_2$}{H$''$}{e$_a$ $\olessthan$ e$_b$}
% }

% \inferrule{I-Combine}{
%   \inline{E}{H}{e$_1$}{H$'$}{e$_a$} \andalso
%   \inline{E}{H$'$}{e$_2$}{H$''$}{e$_b$}
% }{
%   \inline{E}{H}{e$_1$ $\ovee$ e$_2$}{H$''$}{e$_a$ $\ovee$ e$_b$}
% }

% \inferrule{I-Identity}{
%   \inline{E}{H}{e$_1$}{H$'$}{e$_a$} \andalso
%   \inline{E}{H$'$}{e$_2$}{H$''$}{e$_b$}
% }{
%   \inline{E}{H}{e$_1$ == e$_2$}{H$''$}{e$_a$ == e$_b$}
% }

% \infrule[I-Accessor]{
%   \mbox{
%     \evalBig{E}{H}{a}{H$'$}{v} \andalso
%     \lookup{v}{l}{s; return self.l$_f$}}\\
%   \mbox{
%     \enter{E}{H$'$}{C}{I}{v}{E$_m$}{H$_n$}{C}{I} \andalso
%     \inlineS{E$_m$}{H$_n$}{s}{E$'$}{H$''$}{skip}}
% }{\mbox{
%   \inlineS{E}{H}{a.l(e$_1$,$\ldots$,e$_n$)}{E}{H$''$}{v.l$_f$}
% }}

% This last rule restricts what methods used in constraints can
% return. For this step in the semantics, we only allow accessor methods
% in constraints. Note also that the preconditions only allow argument
% expressions that do not add any new constraints.

% This rule shows that the changes to the method environment that happen
% when executing the statements in the method body are dropped, but it
% allows changes to the heap, constraint store, and identity constraint store.
% The field return expression is translated into the caller context.

% Now follow the rules for inlining statements. These just execute the
% statements normally, except for assignments, which are only allowed to
% locals.

% \inferrule{I-AsgnTmp}{
%   \stepH{E}{H}{C}{I}{x := e}{E$'$}{H$'$}{C}{I}
% }{
%   \inlineS{E}{H}{x := e}{E$'$}{H$'$}{skip}
% }

% The above rule just executes the assignment, but does not set up a
% required equality constraint between the LHS and RHS. We are also only
% allowing assignments to local variables right now, because other
% side-effects are disallowed without setting up some sort of constraint
% network. This, too, will be more flexible in the ``full''
% object-oriented language semantics. (In a practical language we found
% it sometimes useful to include a ``backwards compatibility mode''
% which does just that. We also do give no inlining rules for {\tt once}
% and {\tt always} statements, because these are not allowed to be
% nested. Again, a practical language might want to include support for
% nested constraints in a ``backwards compatibility mode.'')

% % Note also that if the assignment execution fails (for example due to
% % conflicting constraints), the inlining gets stuck dynamically.

% All rules essentially just execute the statements.

% \inferrule{I-SeqSkip}{
%   \inlineS{E}{H}{s}{E$'$}{H$'$}{skip}
% }{
%   \inlineS{E}{H}{skip;s}{E$'$}{H$'$}{skip}
% }

% \inferrule{I-SeqStep}{
%   \inlineS{E}{H}{s$_1$}{E$'$}{H$'$}{skip} \andalso
%   \inlineS{E$'$}{H$'$}{s$_1$}{E$''$}{H$''$}{skip} \andalso
% }{
%   \inlineS{E}{H}{s$_1$;s$_2$}{E$''$}{H$''$}{skip}
% }

% \inferrule{I-IfThen}{
%   \evalBig{E}{H}{e}{H$'$}{true} \andalso
%   \inlineS{E}{H$'$}{s$_1$}{E$'$}{H$''$}{skip}
% }{
%   \inlineS{E}{H}{if e then s$_1$ else s$_2$}{E$'$}{H$''$}{skip}
% }

% \inferrule{I-IfElse}{
%   \evalBig{E}{H}{e}{H$'$}{true} \andalso
%   \inlineS{E}{H$'$}{s$_2$}{E$'$}{H$''$}{skip}
% }{
%   \inlineS{E}{H}{if e then s$_1$ else s$_2$}{E$'$}{H$''$}{skip}
% }

% The above two rules show that we are eagerly choosing one branch of
% the if-statement for inlining. Similarly, the next two rules eagerly
% unroll loops.

% \inferrule{I-WhileDo}{
%   \evalBig{E}{H}{e}{H$'$}{true} \andalso
%   \inlineS{E}{H$'$}{s}{E$'$}{H$''$}{s$'$} \andalso
%   \inlineS{E$'$}{H$''$}{while e do s}{E$''$}{H$'''$}{skip}
% }{
%   \inlineS{E}{H}{while e do s}{E$''$}{H$'''$}{skip}
% }

% \inferrule{I-WhileSkip}{
%   \evalBig{E}{H}{e}{H$'$}{false}
% }{
%   \inlineS{E}{H}{while e do s}{E}{H$'$}{skip}
% }


% \judgment{\stepH{E}{H}{C}{I}{s}{E$'$}{H$'$}{C$'$}{I$'$}}

% The rules for normal execution are mostly unchanged.

% \infrule[S-Asgn]{
%   \mbox{
%     \evalH{E}{H}{e}{v}}\\
%   \mbox{
%     \stay{E}{C$_{E_s}$} \andalso
%     \stay{H}{C$_{H_s}$} \andalso
%     \solveH{E$'$}{H$'$}{(I $\wedge$ C$_{E_s}$ $\wedge$ C$_{H_s}$ $\wedge$ L=v)}}\\
%   \mbox{
%     \stay{E$'$}{C$_{{E'}_s}$} \andalso
%     \stay{H$'$}{C$_{{H'}_s}$} \andalso
%     \tcHC{E$'$}{H$'$}{C} \andalso
%     \solveH{E$''$}{H$''$}{(C $\wedge$ C$_{{E'}_s}$ $\wedge$ C$_{{H'}_s}$)}}
% }{\mbox{
%   \stepH{E}{H}{C}{I}{L := e}{E$'$}{H$'$}{C}{I}
% }}

% Assignment changes to work only for l-values, not all accessors.

% \infrule[S-Once]{
%   \mbox{
%     {\tt C$_0$} = $\rho$ {\tt e} \andalso
%     \inline{E}{H$_0$}{e}{H}{e$'$} \andalso
%     {\tt C$_0'$} = $\rho$ {\tt e$'$} \andalso
%     \tcHC{E}{H}{C$_0'$}
%   }\\
%   \mbox{
%     \stay{E}{C$_{E_s}$} \andalso
%     \stay{H}{C$_{H_s}$} \andalso
%     \solveH{E$'$}{H$'$}{$($C $\wedge$ C$_{E_s}$ $\wedge$ C$_{H_s}$ $\wedge$ C$_0'$$)$}}
% }
% {\mbox{\stepH{E}{H$_0$}{C}{I}{once C$_0$}{E$'$}{H$'$}{C}{I}}}

% \infrule[S-Always]{
%   \mbox{
%     {\tt C$_0$} = $\rho$ {\tt e} \andalso
%     \inline{E}{H$_0$}{e}{H}{e$'$} \andalso
%     {\tt C$_0'$} = $\rho$ {\tt e$'$} \andalso
%     \tcHC{E}{H}{C$_0'$}
%   }\\
%   \mbox{
%     \stay{E}{C$_{E_s}$} \andalso
%     \stay{H}{C$_{H_s}$} \andalso
%     {\tt C}$'$ = ({\tt C} $\wedge$ {\tt C}$_0'$) \andalso
%     \solveH{E$'$}{H$'$}{$($C$'$ $\wedge$ {\tt C}$_{E_s}$ $\wedge$ {\tt C}$_{H_s}$$)$}}
% }
% {\mbox{\stepH{E}{H$_0$}{C}{I}{always C$_0$}{E$'$}{H$'$}{C$'$}{I}}}

% \infrule[S-OnceIdentity]{
%   \mbox{
%     \inline{E}{H$_0$}{a$_0$}{H$_0'$}{a$_0'$} \andalso
%     \inline{E}{H$_0'$}{a$_1$}{H}{a$_1'$} \andalso
%     \tcH{E}{H}{a$_0'$}{T} \andalso
%     \tcH{E}{H}{a$_1'$}{T}}\\
%   \mbox{
%     \stay{E}{C$_{E_s}$} \andalso
%     \stay{H}{C$_{H_s}$} \andalso
%     \solveH{E$'$}{H$'$}{$($I $\wedge$ C$_{E_s}$ $\wedge$ C$_{H_s}$ $\wedge$ a$_0'$=a$_1'$$)$}}\\
%   \mbox{
%     \stay{E$'$}{C$_{{E'}_s}$} \andalso
%     \stay{H$'$}{C$_{{H'}_s}$} \andalso
%     \tcHC{E$'$}{H$'$}{C} \andalso
%     \solveH{E$''$}{H$''$}{(C $\wedge$ C$_{{E'}_s}$ $\wedge$ C$_{{H'}_s}$)}}
% }{\mbox{
%     \stepH{E}{H$_0$}{C}{I}{once a$_0$ == a$_1$}{E$''$}{H$''$}{C}{I}
% }}

% \infrule[S-AlwaysIdentity]{
%   \mbox{
%     \inline{E}{H$_0$}{a$_0$}{H$_0'$}{a$_0'$} \andalso
%     \inline{E}{H$_0'$}{a$_1$}{H}{a$_1'$} \andalso
%     \tcH{E}{H}{a$_0'$}{T} \andalso
%     \tcH{E}{H}{a$_1'$}{T}}\\
%   \mbox{
%     \stay{E}{C$_{E_s}$} \andalso
%     \stay{H}{C$_{H_s}$} \andalso
%     {\tt I}$'$ = ({\tt I} $\wedge$ {\tt a}$_0'$={\tt a}$_1'$) \andalso
%     \solveH{E$'$}{H$'$}{$($I$'$ $\wedge$ C$_{E_s}$ $\wedge$ C$_{H_s}$$)$}}\\
%   \mbox{
%     \stay{E$'$}{C$_{{E'}_s}$} \andalso
%     \stay{H$'$}{C$_{{H'}_s}$} \andalso
%     \tcHC{E$'$}{H$'$}{C} \andalso
%     \solveH{E$''$}{H$''$}{(C $\wedge$ C$_{{E'}_s}$ $\wedge$ C$_{{H'}_s}$)}}
% }{\mbox{
%     \stepH{E}{H$_0$}{C}{I}{always a$_0$ == a$_1$}{E$''$}{H$''$}{C}{I$'$}
% }}

% The rules for creating constraints are amended to inline method calls
% (and execute possible side-effects on the heap.)

\section{Babelsberg/Objects}
\label{sec:babelsberg-objects}

In this final language, we add support for mutable objects, classes,
methods, messages, and inheritance, along with object-oriented
constraint definitions (i.e., constraint definitions that can include
method calls).  This language thus includes all the essential features
of actual object constraint languages such as Babelsberg/R,
Babelsberg/JS, and Babelsberg/Squeak.
% For simplicity, we remove field access entirely --- access to an
% objects variables is only possible through message sends.

We can thus now write constraints on the results of message sends,
such as the following constraint on the x coordinate of the center of
a rectangle:
\begin{verbatim}
always r.center().x() = 100;
\end{verbatim}

% I-Compare r.center().x = 100
%   I-Field r.center().x
%     I-MultiWayCall r.center()
%       I-Var pt
%       \ pt
%       H-Value {x: self.x + self.width / 2, self.y + self.height / 2}
%       \ {x: r.x + r.width / 2, r.y + r.height / 2}
%     \ {x: r.x + r.width / 2, r.y + r.height / 2}
%   \ {x: r.x + r.width / 2, r.y + r.height / 2}.x
%   I-Const 100
%   \ 100
% \ {x: r.x + r.width / 2, r.y + r.height / 2}.x = 10

The constraint here is on the result of sending the message
\verb|center| to \verb|r|, then sending the \verb|x| message to that,
and finally sending the \verb|=| message to the result.  Depending on
how the rectangle is stored, the rectangle's center may well be
computed in the \verb|center| method rather than simply being looked
up; and even \verb|x| might be computed rather than being stored, for
example if the point is stored in polar coordinates.  Further,
expressions such as \verb|a+b| are now treated in an object-oriented
fashion, so that this means ``send the message \verb|+| with the
argument \verb|b| to the object \verb|a|,'' with the meaning of
\verb|+| in this expression (and in constraints such as \verb|a+b=c|)
depending on the class of \verb|a|.

In the earlier Babelsbergs presented in this memo, in the semantics we
assume a solver language that is essentially the same as the program
language. Constraints that are sent to the constraint solvers could
thus be in the program language. This changes for objects. As before,
the solver needs to know about uid-records as well as primitive types
and i-records; but it does not know about methods or inheritance. So
we add a translation in this semantics to inline method calls before
passing constraints to the solver.

To accomplish this, we start with a standard model for an
object-oriented language with instances, classes, messages, methods,
and inheritance, and add constraints on top of that.  Some of the
classes and methods in the object-oriented language have corresponding
primitive types and operations which we can pass directly to the
solver. For example, the classes \verb|Integer|, \verb|Float|, and
\verb|Boolean| in the source language might be mapped to the primitive
integers, floats, and booleans. The semantics includes an automatic
boxing and unboxing between instances of these classes and their
primitive equivalents.  (This boxing and unboxing is part of the
formal semantics, but would not necessarily be done in a practical
implementation, where we would want primitive types to be represented
efficiently.)  The \verb|+| method on \verb|Integer| would expand to
use the $+$ operation on a primitive representation of itself, and
similarly for the \verb|+| method for \verb|Float|, the \verb|or|
method for \verb|Boolean|, and so forth. However, in
Babelsberg/Objects, in contrast to prior languages, we may no longer
mix strings and floats, for example. Previously these both typed as
\verb|PrimitiveType|, but now these are full objects that can have
different structure (e.g., a \verb|floatvalue| field to store the
primitive representation on a \verb|Float| object versus a
\verb|stringvalue| field on a \verb|String|). While for these cases,
the structural compatibility checks could be made to work (e.g., by
using a \verb|primvalue| field for these classes), in general there
will be classes in the object-oriented language that don't map to a
primitive type; and even for a class that does have a corresponding
primitive type, that class can have additional methods that aren't
mapped to primitive operations. For example, a \verb|factorial| method
for \verb|Integer| has no primitive equivalent and has to be
translated for the solver.

As with the previous languages, in the semantics, after every
statement execution the current set of constraints is solved and the
environment is updated.  The current set of constraints is determined
as follows.  If the statement is an assignment, we evaluate the right
hand side and then set up a constraint between the resulting value and
the left hand side.  (As with Babelsberg/UID, and in contrast to the
simpler languages, this will be an identity constraint rather than an
equality constraint.)  Otherwise, it must be a method call, which is
evaluated by calling the method normally.
In addition to any constraint resulting from an
assignment statement, each constraint in the constraint store is added
to the set of constraints as well, and the resulting set of
constraints is given to the solver.  If the solver finds a solution,
the environment is updated; but if the constraints have no solution,
or if they are too hard for the solver, our execution implicitly gets stuck.

%% OLD MATERIAL ON CONSTRAINT CONSTRUCTION MODE:
%%
%% Also as before, the evaluation of the closures in the constraint store
%% is done in constraint construction mode rather than in imperative
%% execution mode.  This is more complex than in the previous languages
%% since constraint expressions might invoke methods.  In constraint
%% construction mode, when we encounter a message send that corresponds
%% to an operation on a primitive type in the solver language, we
%% construct the corresponding parse tree fragment for the solver
%% language.  Otherwise, we do a normal method lookup and evaluate the
%% body of the method.  If we we encounter something we cannot transform
%% (such as a side-effect), we don't transform it, and instead execute.
%% As before, assignment statements are converted to \texttt{once}
%% identity constraints with the right-hand-side annotated as read-only.

%% There are two interpretations or meanings \ab{what's the right word?}
%% for each constraint at a particular step in the execution of a program
%% in the Babelsberg/Objects semantics: first, saying the constraint is
%% satisfied is equivalent to saying that the result of evaluting the
%% constraint expression in imperative execution mode is \verb|true|; and
%% second, saying the constraint is satisfied is equivalent to saying
%% that the generated constraint in the solver language is satisfied.
%% Ideally, we would be able to prove that these two interpretations are
%% equivalent for all Babelsberg/Objects programs, based on the rules in
%% the formal semantics.  Further, the behavior of an interpreter or
%% compiler for an actual Babelsberg language should match this
%% interpretation --- the result of evaluating the constraint expression
%% in imperative execution model should give the same result as given by
%% the formal semantics.

%% \ab{The above paragraph is a bit confused but I think it captures an
%%   important idea, namely that there are two meanings for constraints
%%   in the formalism, and they should be provably equivalent.  This
%%   reminds me of work on constraint logic programming, in which one
%%   would show that the operational and the model theoretic semantics
%%   were equivalent, at least under some conditions.}

\subsection{Value Classes}

In the formal semantics, we will use instances of \emph{value classes}
as well as instances of ordinary, full-featured, classes.  The
restrictions on value classes will make it easier to use them with
constraints.  A number of existing languages have proposals to support
forms of value classes, such as Java\footnote{\url{http://openjdk.java.net/jeps/169}} and Scala\footnote{\url{http://docs.scala-lang.org/sips/completed/value-classes.html}}.
In Babelsberg/Objects, instances of value
classes are immutable (that is, we need to provide values for all the
fields at object creation time, and after that the value class instance
cannot be modified).\footnote{For use with e.g.\ distributed systems 
applications, we would also want to restrict instances of value
classes to only hold references to primitive types or to other value
class, but this restriction isn't needed for our purposes here.}
Finally, object identity is not significant for
value class instances --- we define the identity test method, but it
performs a test for equality rather than identity on its argument.
However, value classes are more than simple record declarations, since
they support methods and inheritance.

In this section, we will use \verb|Point| and \verb|Rectangle| as our
example value classes.  To distinguish them from ordinary classes, we
will create instances just by supplying the arguments, e.g.,
\verb|Point(10,20)|, whereas instances of ordinary classes will be
created using the \verb|new| message, e.g., \verb|Window.new(...)|.


In a practical implementation of Babelsberg, ideally the host language
will itself support value classes, so that we can use them directly.
If not, we can use ordinary classes, with appropriate conventions
about object creation, modification, and not testing for object identity.

Finally, to foreshadow the formal semantics, for simplicity we omit
details about object creation, method lookup, and inheritance (which
are completely standard), and represent ordinary instances as
uid-records and value class instances as i-records.

\subsection{Control Structures and Methods}

Babelsberg/Objects has the simple \verb|if| and \verb|while| control
structures that were introduced for Babelsberg/PrimitiveTypes.  There
are no particular complications in including these in
Babelsberg/Objects.

In addition, however, Babelsberg/Objects includes methods.  Following
our overarching design goal of having a standard object-oriented
language in the absence of constraints, in imperative execution mode,
methods are standard.  Suppose we are evaluating \verb|x.m(a,b)|.  We
first look in \verb|x|'s class for a method with selector \verb|m|,
then its superclass, and so on up the superclass chain.  If no such
method is found, it is an error.\footnote{This is analogous to the
  structural compatibility checks from Babelsberg/Records, and just as
  in Babelsberg/Records we don't invent a new field for a record if
  needed to satisfy a constraint, in Babelsberg/Objects we don't
  convert an object to an instance of a different class, or synthesize
  a new method, to satisfy constraints --- instead, if there is a
  missing method it's just an error.  This lookup process isn't
  represented in our formal semantics since it is standard and doesn't
  introduce any issues for our focus on constraints.}
Otherwise we evaluate the body of the method in a new environment, and
return the result.  Arguments are passed by setting up required
identity constraints between the actual and formal parameters.  For
methods that don't involve constraints, this is equivalent to call by
value (with pointer semantics), as in most standard object-oriented
languages.  Again in the absence of constraints, for an argument that
is an instance of a primitive type or value class, it can also be
copied (since we can't tell the difference between sharing and copying
in this case).

However, methods called in the ordinary way can create new
constraints, which then persist after the method returns.  In this
case it does matter that arguments are passed by setting up required
identity constraints between the actual and formal parameters, rather
than using call by value.  There is an example in the next section to
illustrate this.

Methods can also be invoked by constraint expressions.  There are a
number of restrictions on methods to allow them to be used in this way
--- briefly, they can't have side effects, and if they consist of more
than just a single return statement they can only be used in the
forward direction in constraints.  Creating an instance of an ordinary
class is regarded as a side effect, so methods that can be invoked by
constraint expressions can only create instances of value classes.
See Section \ref{sec:additional-constraint-restrictions} for details,
including a discussion of the rationale for this particular choice of
restrictions.

If the method can be used in a constraint but only in the forward
direction, then it can still be called in the same way as with
ordinary method invocations, including passing any parameters by
setting up required identity constraints between the corresponding
actual and formal parameters.  For example,
suppose we have a constraint \verb|always c=x.m(a,b)|, and method
\verb|m| has multiple statements so that it can only be used in the
forward direction, i.e., we can find a value for \verb|c| given values
for \verb|x|, \verb|a|, and \verb|b| but not any other direction.  In
that case, \verb|m| can be called in the usual way, passing \verb|a|
and \verb|b| by value with pointer semantics.  If the other
constraints are such that the correct solution involves finding a
value for \verb|b| given values for \verb|c|, \verb|x|, and \verb|a|,
a practical system will halt with an error that the constraints are
too hard for it to solve, due to a method called from the constraint
that can't be inlined (or in the formal semantics, we get stuck).
However, in contrast to all the previous examples of constraints that
were too hard, the fact that they are too hard is a limitation of the
transformations we use for methods called by constraints, rather than
being a limitation of the solver.

On the other hand, if method \verb|m| can be used multi-directionally,
so that for example we can find a value for \verb|b| given values for
\verb|x|, \verb|a|, and \verb|c|, then in the semantics the method is
inlined so that the resulting constraints can be turned over to the
solver, and we can potentially solve for any of \verb|x|, \verb|a|,
\verb|b|, or \verb|c|.  This done by creating an environment for
method \verb|m| and inlining the returned expression with respect to
that environment. The parameters in that environment are constrained
to be equal to \verb|a| and \verb|b|, and \verb|self| is constrained
to be equal to \verb|x|.

\subsection{Examples}

Here is a constraint on the center of a rectangle, where the center is
a computed value rather than being stored as an instance variable.

\begin{verbatim}
always r.center() = Point(10,20);
\end{verbatim}

The \verb|center| method for \verb|Rectangle| is defined as follows:
\begin{verbatim}
def center
  return (self.upper_left + self.lower_right) / 2;
end
\end{verbatim}

Note that since \verb|Point| is a value class, we can make an instance
of it in the constraint expression itself (``\verb|Point(10,20)|''),
and also in the \verb|center| method, since we are doing point
addition and division in that method.  (If it had been an ordinary
class, creating such an instance would have been a side effect,
disallowing this constraint.)

The constraint is evaluated in the following way.  The code for
Rectangle's \verb|center| method is found, and inlined to construct an
equality constraint between the result of evaluating
\verb|(self.upper_left + self.lower_right) / 2| and a new point with
\verb|x=10|, \verb|y=20|. Constructing the center point from the first
expression, as well as looking up and inlining the \verb|=| method on
it, explodes into a network of simpler constraints that can then be
handed to the solver in one conjunction:
\begin{verbatim}
// for the center method
self1 = r
self2 = self1.upper_left
arg1 = self1.lower_right
argP1_1 = self2.x + arg1.x
argP1_2 = self2.y + arg1.y
self3 = {x: argP1_1, y: argP2_2}
arg2 = 2
argP2_1 = self3.x / arg2
argP2_2 = self3.y / arg2

// for the static point
argP3_1 = 10
argP3_2 = 20

// for the point equality
self4 = {x: argP2_1, y: argP2_2}
arg3 = {x: argP3_1, y: argP3_2}
self4.x = arg3.x && self4.y = arg3.y
\end{verbatim}

For each variable name in a scope a unique global variable name is
created.  The solver works only on the global names. During
evaluation, each local name maps to a global name, which in turn maps
to a value.  Entering the center method creates a local environment in
which \verb|self| is bound to \verb|r|\@.  A global alias \verb|self1|
is created and constrained to refer to the same reference (which is to
a rectangle).  Similarly, the global alias \verb|self2| is a point,
namely the upper left corner of the rectangle.  This alias is then
used when we explode the \verb|+| message to that point.  \verb|self3|
is also a point, namely the new point that is returned by the point
addition method.  Note that we are representing the instance of this
value class in the above conjunction of exploded constraints and in
the formal semantics as an i-record.  For methods that
take arguments, global names for the argument names are created and
constrained similarly.  For example, \verb|arg1| is the global name in
this exploded constraint for the argument to the \verb|+| method for
\verb|Point|. Thus, the solver never has to deal with different scopes
or name clashes.

This example illustrates why we need to solve for object identity
before values --- if \verb|r| had been an instance of a class with a
different implementation of the \verb|center| method, it would have
been exploded in a different way.

% \begin{verbatim}
%   self1 = r // enter center method
%     self2 = self1.upper_left // enter + method
%     arg1 = self.upper_right  // ...
%       argP1_1 = self2.x + arg1.x // enter Point constructor
%       argP1_2 = self2.y + arg1.y // ...
%       ^ {x: argP1_1, y: argP2_2}
%     self3 = {x: argP1_1, y: argP2_2} // enter / method
%     arg2 = 2                         // ...
%       argP2_1 = self3.x / arg2 // enter Point constructor
%       argP2_2 = self3.y / arg2 // ...
%       ^ {x: argP2_1, y: argP2_2}
%     ^ {x: argP2_1, y: argP2_2}
%   ^ {x: argP2_1, y: argP2_2}

%   argP3_1 = 10 // enter Point constructor
%   argP3_2 = 20 // ...
%   ^ {x: argP3_1, y: argP3_2}

%   self4 = {x: argP2_1, y: argP2_2} // enter = method
%   arg3 = {x: argP3_1, y: argP3_2}  // ...
%   ^ self4.x = arg3.x && self4.y = arg3.y
% \end{verbatim}

Here is another example of using a method in a constraint.  Suppose we
add a \verb|double| method to \verb|Float|, and then use it in a
constraint:

\begin{verbatim}
def double()
  return 2*self;
end;

x := 0;
y := 0;
always y=x.double();
y := 20;
\end{verbatim}

After the program runs, \verb|y| will be 20 and \verb|x| will be 10.

For this program, the constraint \verb|y=x.double()| explodes into the
following network:

\begin{verbatim}
self1 = x
y = 2*self1
\end{verbatim}

This conjunction of constraints works equally well for determining
\verb|x| or determining \verb|y|, even though \verb|double| is
actually a message to \verb|x|.

As noted above, methods called in the ordinary way can create new
constraints, which then persist after the method returns.  Here is a
simple example.  Suppose we have a \verb|BankAccount| class that
includes a \verb|balance| method.  We can then define the following method:

\begin{verbatim}
def require_min_balance(acct,min)
  always acct.balance >= min;
end;
\end{verbatim}

Now suppose we call \verb|require_min_balance(a,10)| on some account
\verb|a|.  (This is just calling the method from an ordinary
statement, either at the top level or some other method, \emph{not}
from a constraint expression.)  Thereafter the balance in the account
\verb|a| must be at least 10 (a persistent constraint).

Note the consequences of call by value for such methods.  Consider:
\begin{verbatim}
a := BankAccount.new;
m := 10;
require_min_balance(a,m);
m := 100;
\end{verbatim}
Even though we reassigned \verb|m|, the minimum balance for the
account stays at 10, since we passed \verb|m| by value.

Of course, we are passing the account \verb|a| by value as well.

\begin{verbatim}
a := BankAccount.new;
require_min_balance(a,10);
a := BankAccount.new;
\end{verbatim}

After the second assignment,  \verb|a| is bound to a new account --- but
the minimum balance constraint is on the previous instance of
\verb|BankAccount| (which will presumably be garbage collected since
there aren't any references to it).  

If we do want a constraint that continues to be enforced even if the
account or the minimum is rebound, we can instead write a method that
is a minimum balance test and use it in a constraint:

\begin{verbatim}
def has_min_balance(acct,min)
  return acct.balance >= min;
end;

a := BankAccount.new;
m := 10
always has_min_balance(a,m);
m := 100;
\end{verbatim}

Programmers will need to be aware of the different semantics for these
two cases, and select the appropriate variant when it makes a difference.

Implementing such methods requires that the semantics (and
implementations) keep the local variables in such a method invocation
as long as the constraints it creates are active.  Here is a somewhat
artificial variant of the original \verb|require_min_balance|
method to illustrate this:
\begin{verbatim}
def require_min_balance_with_locals(acct,min)
  b := acct.balance;
  always b = acct.balance;
  always b >= min;
end;
\end{verbatim}

We could also add a minimum balance constraint in the initialization
method for \verb|BankAccount|.  This version requires that the initial
balance be provided, but has a default of 0 for the minimum balance in
case it isn't provided.

\begin{verbatim}
def init(initial_balance,min=0)
  self.balance := initial_balance;
  always self.balance >= min;
end;
\end{verbatim}

\subsection{Arrays}

While it is not strictly necessary to include arrays in
Babelsberg/Objects, we have used arrays in some examples that follow.
We treat the length of an array is part of its structure, and use the
same structural compatibility checks for the existence of array
elements that we do for the existence of fields in a record.  Thus, we
would not, for example, grow an array to allow a constraint on its
$i^{th}$ element to be satisfied if it didn't already have an $i^{th}$
element.

\subsection{Additional Restrictions on Constraint Expressions}
\label{sec:additional-constraint-restrictions}

Methods can of course have side effects, but 
such methods can't be used in the expressions that define
constraints.  (This is one of the restrictions noted in Section
\ref{sec:constraint-expression-requirements} on requirements for
constraint expressions.)  This subsection describes some consequences
of this restriction when we have objects, methods, and object
identity.  There are also some additional restrictions on methods
called from constraints that are used in other than the forward
direction, discussed below.

Regarding side effects in methods, consider the \verb|pop| method for
a class \verb|Stack|, which has a side effect as well as returning the
value popped from the stack.  This program fragment is OK:

\begin{verbatim}
x := stack.pop;
\end{verbatim}

This works, because the RHS of the assignment is evaluated first, and
then we set up a \verb|once| constraint using the value popped from the
stack.  (Thus the RHS of an assignment can have side effects, in contrast
to \verb|always| or \verb|once| constraints.)

In contrast, these program fragments are not OK:
\begin{verbatim}
once x == stack.pop;
\end{verbatim}

\begin{verbatim}
always x == stack.pop;
\end{verbatim}

It's easy to see why we don't want to allow this: we should be able to
evaluate the constraint expression as a test, and if it evaluates to
true the constraint is satisfied.  This would just not work in these
stack examples.

As noted previously, creating an instance of an ordinary class is
regarded as a side effect, so methods that can be invoked by
constraint expressions can only create instances of value classes and
not ordinary classes.\footnote{As also noted previously (Footnote
  \ref{footnote:benign-side-effects}), in a practical implementation
  the programmer might be able to make cautious use of benign side effects
  in a constraint expression.  Such benign side effects might include
  creating instances of ordinary classes, for example, constructing a
  temporary instance of an ordinary class that is garbage collected
  before it is visible outside the constraint.  In the formal
  semantics, however, we simply disallow side effects in constraint
  expressions.}

Creating a constraint is a kind of side effect, at least potentially,
so another consequence of the prohibition on side effects in
constraint expressions is that methods called from a constraint
expression cannot themselves create other constraints, or call further
methods that do so.  The following program is not allowed, and
illustrates side effects resulting from adding constraints in methods
called from other constraints.

\begin{verbatim}
def test(i)
  always medium i=5;
  return i+1;
end;
x := 0;
y := 0;
always medium x=10;
always y=test(x);
\end{verbatim}

Suppose we are using a least-squares solver that supports soft
constraints.  Every time we evaluate the expression \verb|y=test(x)|,
we add another medium constraint that \verb|x=5|.  Since we are
finding a least squares solution, this nudges \verb|x| more toward 5
each time, diminishing the influence of the \verb|x=10|
constraint.\footnote{This issue only arises with soft constraints and
  a global comparator.  With constraints that are required, re-adding
  a required constraint is idempotent.  Or if we use a local
  comparator, re-adding the soft constraint will also not change the
  resulting solution.  But in any case, adding a constraint is
  conceptually different from testing whether it is satisfied.}

Here are two other issues with constraints that call methods that add
constraints, for the record in case we do want to try to allow this in
some restricted set of circumstances:

First, we would need to solve their constraints eagerly rather than
just accumulating the constraints and solving all at once.  Here's an
example that demonstrates this.

\begin{verbatim}
def test(x)
  always y=x;
  if y>10 then return 1 else return 0;
end;
a := 5;
always b=test(a);
\end{verbatim}
If we don't solve eagerly, \verb|y| doesn't have a value when we
evaluate the test in the \verb|if| statement.

Also the method could add soft constraints, and itself be invoked by a
soft constraint --- for example, suppose that in the above program we
instead had \verb|always weak b=test(a)|.  We would need an algebra
for combining priorities, for example, take the minimum of the
priorities.

Methods called by constraints can assign to local variables, but not
to instance variables or globals (which would be side effects visible
outside the method).  If there are assignments to local variables, the
method can only be used in the forward direction in constraints that
call it (i.e., to compute the result given the inputs).  The only way
that statements in a method other than the final return aren't dead
code is if they have side effects, so for practical purposes we can
restate this restriction as follows: for a method to be used in other
than just the forward direction in constraints that call it, the
method must consist of a single \verb|return| statement.  

It would be possible to relax this restriction somewhat, so that other
kinds of methods could be used in reverse in constraints, but the
resulting restrictions would be harder to express, and would have an
\emph{ad hoc} feel, likely making it harder for the programmer to keep
track of them.  So in the design presented here we use a simple,
easily understood restriction that covers an important set of
practical cases.  See Appendix~\ref{sec:backwards-compatibility-mode}
for a description of a backward compatibility mode for methods that
accommodates some methods with multiple statements, which is currently
implemented in Babelsberg/Ruby.

Here are some examples.

Consider an iterative \verb|sum| method for
the class \verb|Array|.

\begin{verbatim}
def sum
  ans := 0;
  i := 0;
  while i<self.length
    ans := ans + self[i];
    i := i+1;
  end;
  return ans;
end
\end{verbatim}

This works as expected for normal imperative code, and can also be
used in the forward direction in a constraint, for example:

\begin{verbatim}
a := Array.new(2);
a[0] := 10;
a[1] := 20;
s := 30;
always s = a.sum;
a[0] := 100;
\end{verbatim}

After the \verb|always| constraint is executed, \verb|s| is 30; then
after the final assignment to \verb|a[0]|, \verb|s| becomes 120.

However, the method doesn't work backwards --- for example, we can't
constrain the sum of the array and expect the system to update one or
more elements to satisfy the constraint.  So the constraint in the
last line below will be too hard for the system to solve:

\begin{verbatim}
a := Array.new(2);
a[0] := 10;
a[1] := 20;
always 50 = a.sum;
\end{verbatim}

We \emph{can} provide something analogous to an array \verb|sum|
constraint that works both forward and backward by writing an ordinary
method that sets up the appropriate network of addition constraints.
It is defined here as the \verb|sum_equal| method on \verb|Array|,
although it could be a standalone method instead.

\begin{verbatim}
def sum_equal(sum)
  helper(0,sum);
end;

def helper(start,sum)
  if start>=self.length then always sum=0;
    else 
       partial := 0;  /* create a local variable */
       helper(start+1,partial);
       always sum=self[start]+partial;
  end if;
end;  
\end{verbatim}

Now our example works:

\begin{verbatim}
a := Array.new(2);
a[0] := 10;
a[1] := 20;
a.sum_equal(50);
\end{verbatim}

If we subsequently set \verb|a[1]| to 5, \verb|a[0]| will become 45 to
keep the constraints satisfied.

Note that in contrast to the earlier version, the constraints are not
automatically re-applied if \verb|a| is reassigned; the programmer
needs to call \verb|sum_equal| again.\footnote{At the cost of
  complicating the language, we could address this issue by adding a
  construct that automatically invokes \texttt{some\_method} whenever
  needed to re-establish the constraint network.  To be used in this
  way, \texttt{some\_method} could only assert constraints; it could
  have no other side effects.  There'd also need to be machinery to
  retract the old constraints before re-running \texttt{some\_method},
  to avoid accumulating obsolete constraints.}

It is tempting to write \verb|sum| as a method that returns the sum
and that still can be used multi-directionally, perhaps:

\begin{verbatim}
def sum
  return self.sum_from(0);
end;

def sum_from(start)
  if start >= self.length
    then return 0
    else return self[start] + self.sum_from(start+1);
  end;
end;
\end{verbatim}

Unfortunately, this doesn't meet the restriction that methods called
from constraints consist only of a single return statement.  Worse, if
we tried inlining it, the method would expand infinitely.  So, at
least for now, \verb|sum_equal| seems like it is as good as we can do.

A related example is a method for \verb|Array| that tests whether the
argument is one of the elements of that array.

\begin{verbatim}
def contains(x)
  i:=0;
  while i<self.length && self[i]!=x
    i:=i+1;
  end;
  return i<self.length;
end;
\end{verbatim}

Since this includes more than one statement, it can only be used in
a constraint in the forward direction:
\begin{verbatim}
a := Array.new(10);
a.fill(100);
c := false;
always c = a.contains(5);
a[3] := 5;
\end{verbatim}

The analogous method to \verb|sum_equal|, say \verb|must_contain|,
unfortunately would not be very useful: since we don't have
Prolog-style backtracking, it would probably just be satisfied in the
reverse direction by setting the first element of the array.

\subsection{Identity Constraint Examples}

The following example illustrates how assignment for objects with
object identity is handled.  This uses an example class
\verb|MutablePoint|, which is an ordinary class whose instances can be
changed and that have individual identities, in contrast to the value
class \verb|Point| used previously.

\begin{verbatim}
p1 := MutablePoint.new(10,10);
p2 := p1;
p1 := MutablePoint.new(50,50);
\end{verbatim}

The result is just the same as in Ruby or some other object-oriented
language without constraints.  After the first statement, we have the
new point with \verb|x=10,y=10| in the heap, and solve a required identity
constraint that \verb|p1|'s value equal the reference to this point.
After the second statement, the identity constraints are a weak
identity stay on \verb|p1|, and a required constraint that \verb|p2|
reference the same value as \verb|p1|.  After the third statement,
there is another point with \verb|x=50,y=50| in the heap, a required
constraint that \verb|p1|'s value now equal the reference to this new
point, and weak identity constraints that \verb|p1| refer to what it
used to (which can't be satisfied), and that \verb|p2| refer to what
it used to (which \emph{can} be satisfied.)

As with Babelsberg/UID, in the interests of predictability, the system
will not create a new object to satisfy constraints.  Thus, this
program halts with an error:
\begin{verbatim}
p := MutablePoint.new(0,0);
q := p;
always p.x=5;
always q.x=10;
\end{verbatim}

After the third statement, both \verb|p| and \verb|q| refer to the
same point, which has \verb|x=5,y=0|.  The final constraint on 
\verb|q.x| is unsatisfiable.  We could satisfy the constraints by
first cloning \verb|q| from \verb|p| and then changing its \verb|x|,
but we forbid creating new objects to satisfy constraints.  In the
semantics, this results in getting stuck due to an unsatisfiable
required constraint.

While cloning a mutable point seems innocuous, consider a similar example with
windows on the user's display, which again results in an error.  (It
seems dangerous to silently create a second window.)

\begin{verbatim}
w1 := Window.new(....);
w2 := w1;
always w1.width = 200;
always w2.width = 300;
\end{verbatim}

However, this variant, in which we make the constraints on the
window's widths strong rather than required, is OK.  (We don't create
a new window to satisfy the strong constraints on the widths.)

\begin{verbatim}
w1 := Window.new(....);
w2 := w1;
always strong w1.width = 200;
always strong w2.width = 300;
\end{verbatim}

As a point of comparison, this restriction on not creating new objects
isn't relevant for Babelsberg/Records or when we use instances of
value classes in Babelsberg/Objects.
Here is one of the above examples using value classes:

\begin{verbatim}
p := Point(10,20);
q := p;
always p.x=5;
always q.x=10;
\end{verbatim}

This is fine, and at the end of the program, \verb|p| and \verb|q|
will refer to two different value class instances.

There can be explicit identity constraints, as in Babelsberg/UID\@.
Again as in Babelsberg/UID, identity constraints must be satisfied
when they are first created.  For example:
\begin{verbatim}
x:= Window.new(....);
y := x;
always y==x;
x := Circle.new(....);
\end{verbatim}


Here \verb|x| and \verb|y| are identical at the time the identity
constraint \verb|always y==x| is created.  And given this identity
constraint, after the last assignment both \verb|x| and \verb|y| refer
to the newly created circle.  This example also illustrates how type
changes must originate with an assignment statement, but can ripple
further via explicit identity constraints.

\subsection{Formalism}

\subsubsection{Syntax}

The syntax is augmented to support method invocation as expressions. We introduce
syntax for the method body.
Additionally, we allow creating new objects as expressions
now. We also have a form of immutable records and consider them as
``value objects,'' which can respond to methods.   Finally, we
introduce the {\tt nil} value.  The entire syntax is given below.

\begin{tt}
\begin{tabular}{lrrl}
{\rm Statement} & s & ::= & skip $\mid$ L := e $\mid$ always C $\mid$ once C $\mid$ s;s \\
  & & & $\mid$ if e then s else s $\mid$ while e do s \\
{\rm Constraint} & C & ::= & $\rho$ e $\mid$ C $\wedge$ C\\
{\rm Expression} & e & ::= & v $\mid$ L $\mid$ e $\oplus$ e
$\mid$ e $\olessthan$ e $\mid$ e $\ovee$ e $\mid$ I $\mid$
e.l(e$_1$,$\ldots$,e$_n$) $\mid$ o $\mid$ new o $\mid$ D\\
{\rm Identity} & I & ::= & e == e\\
{\rm Object Literal} & o & ::= &  \{l$_1$:e$_1$,$\ldots$,l$_n$:e$_n$\} \\
{\rm L-Value} & L & ::= & x $\mid$ e.l \\
{\rm Constant} & c & ::= & true $\mid$ false $\mid$ nil $\mid$ {\rm base type constants}\\
{\rm Variable} & x & ::= & {\rm variable names} \\
{\rm Label} & l & ::= & {\rm record label names} \\
{\rm Reference} & r & ::= & {\rm references to heap records} \\
{\rm Dereference} & D & ::= & \H(e) \\\\
{\rm Method Body} & b & ::= & s; return e $\mid$ return e \\

{\rm Value} & v & ::= & c $\mid$ r $\mid$ \{l$_1$:v$_1$,$\ldots$,l$_n$:v$_n$\} \\
\end{tabular}
\end{tt}

We assume that the set of variable names ranged over by metavariable {\tt
  x} includes the name {\tt self}.

\subsubsection{Semantics}

Since we have methods with local scopes now, we introduce a global \E\
that maps global variable names to values. The rules for execution
still work on a local environment, but a local environment is only a
mapping from local names to globally unique names now. Method lookup
creates new local environments, and constraints are translated to
translate local variable names to global variable names for the
solver. The solver still returns a global environment, and does not
know about the local environments and their mapping into the global
environment. In addition, the constraint stores {\tt I} and {\tt C}
are replaced with \I\ and \C, respectively. These now store
constraints in pairs with the local environment they were created
in. For readability, we now write \H\ instead of {\tt H}, so all the
global stores are written in the same font.

\judgment{\lookup{v}{l}{b}}

``Lookup of method {\tt l} in the object or value {\tt v} returns the
formal parameter names {\tt x$_1$} through {\tt x$_n$} and the method
body {\tt b}''

This judgment is opaque:  our semantics does not depend on how method
lookup is performed.

\judgment{\enter{\E}{E}{\H}{\C}{\I}{v}{\E$'$}{E$_m$}{\H$'$}{\C$'$}{\I$'$}}

``Invoking a method on {\tt v} with argument names {\tt x$_1$} through
{\tt x$_n$} and arguments {\tt e$_1$} through {\tt e$_n$} constructs
the method scope {\tt E$_m$} and may update the heap and constraint stores.''

This is a helper judgment that simplifies the definition of evaluation
for method calls (shown below).

\infrule[Enter]{
  \mbox{
    \evalBig{\E}{E}{\H}{\C}{\I}{e$_1$}{\E$_1$}{\H$_1$}{\C$_1$}{\I$_1$}{v$_1$}}\\
  \mbox{$\cdots$}\\
  \mbox{
    \evalBig{\E$_{n-1}$}{E}{\H$_{n-1}$}{\C$_{n-1}$}{\I$_{n-1}$}{e$_n$}{\E$_n$}{\H$_n$}{\C$_n$}{\I$_n$}{v$_n$}}\\
  \mbox{
    \stepHO{\E$_n$}{E$_m$}{\H$_n$}{\C$_n$}{\I$_n$}{self := v}{\E$_0$}{E$_0$}{\H$_n$}{\C$_n$}{\I$_n$}}\\
  \mbox{
    \stepHO{\E$_0$}{E$_0$}{\H$_n$}{\C$_n$}{\I$_n$}{x$_1$ := v$_1$}{\E$_{n+1}$}{E$_1$}{\H$_n$}{\C$_n$}{\I$_n$}}\\
  \mbox{$\cdots$}\\
  \mbox{
    \stepHO{\E$_{2n-1}$}{E$_{n-1}$}{\H$_n$}{\C$_n$}{\I$_n$}{x$_n$ := v$_n$}{\E$_{2n}$}{E$_n$}{\H$_n$}{\C$_n$}{\I$_n$}}
}{\mbox{
    \enter{\E}{E}{\H}{\C}{\I}{v}{\E$_{2n}$}{E$_n$}{\H$_n$}{\C$_n$}{\I$_n$}
}}

\judgment{\evalBig{\E}{E}{\H}{\C}{\I}{e}{\E$'$}{\H$'$}{\C$'$}{\I$'$}{v}}

``Evaluating expression {\tt e} produces the value {\tt v}, while
possibly having side-effects on everything but the local environment.''

\inferax{E-Const}{\evalBig{\E}{E}{\H}{\C}{\I}{c}{\E}{\H}{\C}{\I}{c}}

\inferrule{E-Var}{
  {\tt E}({\tt x}) = {\tt x$_g$} \andalso \E({\tt x$_g$}) = {\tt v}
}{
  \evalBig{\E}{E}{\H}{\C}{\I}{x}{\E}{\H}{\C}{\I}{v}
}

\inferrule{E-Field}{
  \evalBig{\E}{E}{\H}{\C}{\I}{e}{\E$'$}{\H$'$}{\C$'$}{\I$'$}{r} \andalso
  {\tt \H$'$}({\tt r}) = {\tt \{l$_1$:v$_1$,$\ldots$,l$_n$:v$_n$\}} \andalso
  1 $\leq$ $i$ $\leq$ $n$
}{
  \evalBig{\E}{E}{\H}{\C}{\I}{e.l$_i$}{\E$'$}{\H$'$}{\C$'$}{\I$'$}{v$_i$}
}

\inferrule{E-ValueField}{
  \evalBig{\E}{E}{\H}{\C}{\I}{e}{\E$'$}{\H$'$}{\C$'$}{\I$'$}{\{l$_1$:v$_1$,$\ldots$,l$_n$:v$_n$\}} \andalso
  1 $\leq$ $i$ $\leq$ $n$
}{
  \evalBig{\E}{E}{\H}{\C}{\I}{e.l$_i$}{\E$'$}{\H$'$}{\C$'$}{\I$'$}{v$_i$}
}

\inferax{E-Ref}{\evalBig{\E}{E}{\H}{\C}{\I}{r}{\E}{\H}{\C}{\I}{r}}

\inferrule{E-Op}{
  \evalBig{\E}{E}{\H}{\C}{\I}{e$_1$}{\E$_0$}{\H$_0$}{\C$_0$}{\I$_0$}{v$_1$} \andalso
  \evalBig{\E$_0$}{E}{\H$_0$}{\C$_0$}{\I$_0$}{e$_2$}{\E$'$}{\H$'$}{\C$'$}{\I$'$}{v$_2$} \andalso
  {\tt v$_1$ $\llbracket \oplus \rrbracket$ v$_2$} = {\tt v}
}{
  \evalBig{\E}{E}{\H}{\C}{\I}{e$_1$ $\oplus$ e$_2$}{\E$'$}{\H$'$}{\C$'$}{\I$'$}{v}
}

\inferrule{E-Compare}{
  \evalBig{\E}{E}{\H}{\C}{\I}{e$_1$}{\E$_0$}{\H$_0$}{\C$_0$}{\I$_0$}{v$_1$} \andalso
  \evalBig{\E$_0$}{E}{\H$_0$}{\C$_0$}{\I$_0$}{e$_2$}{\E$'$}{\H$'$}{\C$'$}{\I$'$}{v$_2$} \andalso
  {\tt v$_1$ $\llbracket \olessthan \rrbracket$ v$_2$} = {\tt v}
}{
  \evalBig{\E}{E}{\H}{\C}{\I}{e$_1$ $\olessthan$ e$_2$}{\E$'$}{\H$'$}{\C$'$}{\I$'$}{v}
}

\inferrule{E-Combine}{
  \evalBig{\E}{E}{\H}{\C}{\I}{e$_1$}{\E$_0$}{\H$_0$}{\C$_0$}{\I$_0$}{v$_1$} \andalso
  \evalBig{\E$_0$}{E}{\H$_0$}{\C$_0$}{\I$_0$}{e$_2$}{\E$'$}{\H$'$}{\C$'$}{\I$'$}{v$_2$} \andalso
  {\tt v$_1$ $\llbracket \ovee \rrbracket$ v$_2$} = {\tt v}
}{
  \evalBig{\E}{E}{\H}{\C}{\I}{e$_1$ $\ovee$ e$_2$}{\E$'$}{\H$'$}{\C$'$}{\I$'$}{v}
}


\infrule[E-IdentityTrue]{
  \mbox{%
  \evalBig{\E}{E}{\H}{\C}{\I}{e$_1$}{\E$_0$}{\H$_0$}{\C$_0$}{\I$_0$}{v} \andalso
  \evalBig{\E}{E}{\H$_0$}{\C$_0$}{\I$_0$}{e$_2$}{E$'$}{\H$'$}{\C$'$}{\I$'$}{v}}
}{\mbox{%
  \evalBig{\E}{E}{\H}{\C}{\I}{e$_1$ == e$_2$}{\E$'$}{\H$'$}{\C$'$}{\I$'$}{true}
}}

\infrule[E-IdentityFalse]{
  \mbox{%
  \evalBig{\E}{E}{\H}{\C}{\I}{e$_1$}{\E$_0$}{\H$_0$}{\C$_0$}{\I$_0$}{v$_1$} \andalso
  \evalBig{\E}{E}{\H$_0$}{\C$_0$}{\I$_0$}{e$_2$}{E$'$}{\H$'$}{\C$'$}{\I$'$}{v$_2$}
  \andalso {\tt v$_1$} $\neq$ {\tt v$_2$}}
}{\mbox{%
  \evalBig{\E}{E}{\H}{\C}{\I}{e$_1$ == e$_2$}{\E$'$}{\H$'$}{\C$'$}{\I$'$}{false}
}}

\infrule[E-Call]{\mbox{
  \evalBig{\E}{E}{\H}{\C}{\I}{e}{\E$_0$}{\H$_0$}{\C$_0$}{\I$_0$}{v} \andalso
  \lookup{v}{l}{s; return e}}\\
\mbox{
  \enter{\E$_0$}{E}{\H$_0$}{\C$_0$}{\I$_0$}{v}{\E$_1$}{E$_m$}{\H$_1$}{\C$_1$}{\I$_1$}}\\
\mbox{\stepHO{\E$_1$}{E$_m$}{\H$_1$}{\C$_1$}{\I$_1$}{s}{\E$'$}{E$'$}{\H$'$}{\C$'$}{\I$'$} \andalso
  \evalBig{\E$'$}{E$'$}{\H$'$}{\C$'$}{\I$'$}{e}{\E$''$}{\H$''$}{\C$''$}{\I$''$}{v$_r$}}
}{\mbox{
  \evalBig{\E}{E}{\H}{\C}{\I}{e.l(e$_1$,$\ldots$,e$_n$)}{\E$''$}{\H$''$}{\C$''$}{\I$''$}{v$_r$}
}}

\infrule[E-CallSimple]{\mbox{
  \evalBig{\E}{E}{\H}{\C}{\I}{e}{\E$_0$}{\H$_0$}{\C$_0$}{\I$_0$}{v} \andalso
  \lookup{v}{l}{return e}}\\
\mbox{
  \enter{\E$_0$}{E}{\H$_0$}{\C$_0$}{\I$_0$}{v}{\E$_1$}{E$_m$}{\H$_1$}{\C$_1$}{\I$_1$}}\\
\mbox{
  \evalBig{\E$_1$}{E$_m$}{\H$_1$}{\C$_1$}{\I$_1$}{e}{\E$'$}{\H$'$}{\C$'$}{\I$'$}{v$_r$}}
}{\mbox{
  \evalBig{\E}{E}{\H}{\C}{\I}{e.l(e$_1$,$\ldots$,e$_n$)}{\E$'$}{\H$'$}{\C$'$}{\I$'$}{v$_r$}
}}

\infrule[E-New]{\mbox{
  \evalBig{\E}{E}{\H}{\C}{\I}{e$_1$}{\E$_1$}{\H$_1$}{\C$_1$}{\I$_1$}{v$_1$} $\cdots$ 
  \evalBig{\E$_{n-1}$}{E}{\H$_{n-1}$}{\C$_{n-1}$}{\I$_{n-1}$}{e$_n$}{\E$_n$}{\H$_n$}{\C$_n$}{\I$_n$}{v$_n$}}\\
\mbox{
{\tt \H$_n$}({\tt r})$\uparrow$ \andalso  {\tt \H$'$} = ({\tt \H$_n$} $\bigcup$ \{({\tt r}, {\tt \{l$_1$:v$_1$,$\ldots$,l$_n$:v$_n$\}})\})}
}{\mbox{
  \evalBig{\E}{E}{\H}{\C}{\I}{new \{l$_1$:e$_1$,$\ldots$,l$_n$:e$_n$\}}{\E$_n$}{\H$'$}{\C$_n$}{\I$_n$}{r}
}}

\infrule[E-Value]{\mbox{
  \evalBig{\E}{E}{\H}{\C}{\I}{e$_1$}{\E$_1$}{\H$_1$}{\C$_1$}{\I$_1$}{v$_1$} $\cdots$ 
  \evalBig{\E$_{n-1}$}{E}{\H$_{n-1}$}{\C$_{n-1}$}{\I$_{n-1}$}{e$_n$}{\E$_n$}{\H$_n$}{\C$_n$}{\I$_n$}{v$_n$}}\\
}{\mbox{
  \evalBig{\E}{E}{\H}{\C}{\I}{\{l$_1$:e$_1$,$\ldots$,l$_n$:e$_n$\}}{\E$_n$}{\H$_n$}{\C$_n$}{\I$_n$}{\{l$_1$:v$_1$,$\ldots$,l$_n$:v$_n$\}}
}}

\judgment{\tcH{\E}{\H}{e}{T}}
\judgment{\tcHC{\E}{\H}{C}}

We typecheck with respect to the global environments \E$\,$.  Method
calls do not typecheck: even though we allow method calls in
expressions and thus in constraints syntactically, our rules for
creating constraints given below inline method invocations.  Identity
constraints also do not typecheck: they are solved separately and
should not appear in ordinary constraints (see the rules for statement
evaluation below).  Finally, new (non-value) object construction does
not typecheck, since constraints must be side-effect-free.

Types distinguish between primitive values and objects, and for
objects the type keeps track of their fields:

\begin{tt}
\begin{tabular}{lrrl}
{\rm Type} & T & ::= & PrimitiveType $\mid$
\{l$_1$:T$_1$,$\ldots$,l$_n$:T$_n$\} 
\end{tabular}
\end{tt}

\inferax{T-Constant}{
  \tcH{\E}{\H}{c}{PrimitiveType}
}

\inferrule{T-Variable}{
  \E({\tt x}) = {\tt v} \andalso
  \tcH{\E}{\H}{v}{T}
}{
  \tcH{\E}{\H}{x}{T}
}

We now transform all local variables to their global names using the
inlining rules given below before passing them to the solver. Thus,
only global variable names type.

\inferrule{T-Field}{
  \tcH{\E}{\H}{e}{\{l$_1$:T$_1$,$\ldots$,l$_n$:T$_n$\}}
  \andalso $1 \leq i \leq n$
}{
  \tcH{\E}{\H}{e.l$_i$}{T$_i$}
}

\inferrule{T-Ref}{
  {\tt \H}({\tt r})={\tt o} \andalso
  \tcH{\E}{\H}{o}{T} 
}{
  \tcH{\E}{\H}{r}{T}
}

\inferrule{T-Deref}{
  \tcH{\E}{\H}{e}{T}
}{
  \tcH{\E}{\H}{\H(e)}{T}
}

\inferrule{T-Op}{
  \tcH{\E}{\H}{e$_1$}{PrimitiveType} \andalso \tcH{\E}{\H}{e$_2$}{PrimitiveType}
}{
  \tcH{\E}{\H}{e$_1$ $\oplus$ e$_2$}{PrimitiveType}
}

\inferrule{T-Compare}{
  \tcH{\E}{\H}{e$_1$}{PrimitiveType} \andalso \tcH{\E}{\H}{e$_2$}{PrimitiveType}
}{
  \tcH{\E}{\H}{e$_1$ $\olessthan$ e$_2$}{PrimitiveType}
}

\inferrule{T-Combine}{
  \tcH{\E}{\H}{e$_1$}{PrimitiveType} \andalso \tcH{\E}{\H}{e$_2$}{PrimitiveType}
}{
  \tcH{\E}{\H}{e$_1$ $\ovee$ e$_2$}{PrimitiveType}
}

\inferrule{T-ValueObject}{
  \tcH{\E}{\H}{e$_1$}{T$_1$} $\cdots$ \tcH{\E}{\H}{e$_n$}{T$_n$}
}{
  \tcH{\E}{\H}{\{l$_1$:e$_1$,$\ldots$,l$_n$:e$_n$\}}{\{l$_1$:T$_1$,$\ldots$,l$_n$:T$_n$\}}
}



\inferrule{T-Constraint}{
  \tcH{\E}{\H}{e}{T}
}{
  {\tcHC{\E}{\H}{$\rho$ e}}
}

\inferrule{T-CompoundConstraint}{
  \tcH{\E}{\H}{C$_1$}{T} \andalso \tcH{\E}{\H}{C$_2$}{T}
}{
  \tcHC{\E}{\H}{C$_1$ $\wedge$ C$_2$}
}


\judgment{\solveH{\E}{\H}{C}}

As before we assume that the solver natively
supports records and uninterpreted functions (which we use to
represent the heap).  We do not
assume, however, that the solver understands methods, which can now be
part of constraint expressions. This requires us to essentially inline
methods before passing constraints to the solver.

\judgment{\stay{\E}{C}}
\judgment{\stay{\H}{C}}
\judgment{\istay{e}{C}}
\judgment{\istay{\E}{C}}
\judgment{\istay{\H}{C}}

The rules for creating stay constraints are adapted to work with
multiple environments. Also, there are now two groups of rules. The
first two stay rules assert weak stays as before, but the
\istay{\E}{C} and \istay{\H}{C} rules use the \istay{e}{C} rule to
create weak stays on primitive values and irecords, but required stays
for references. This ensures that value constraints cannot affect
identities.

\inferax{StayEmpty}{\stay{$\emptyset$}{true}}

\inferrule{StayOne}{
  \E$({\tt x}) = {\tt v}$ \andalso
  \E$_0 = $\E$ \setminus \{({\tt x},{\tt v})\}$ \andalso
  \stay{\E$_0$}{C$_0$} \andalso
  {\tt C}$ = ${\tt C}$_0$ $\wedge$ {\tt weak x=v}
}{
  \stay{\E}{C}
}

\inferrule{StayHeap}{
  \H$({\tt r}) = {\tt o}$ \andalso  
  \H$_0 = $\H$ \setminus \{({\tt r},{\tt o})\}$ \andalso
  \stay{\H$_0$}{C$_0$} \andalso
  ${\tt o} = $\{{\tt l$_0$:v$_0$,$\ldots$,l$_n$:v$_n$}\}
}{
  \stay{\H}{C$_0$ $\wedge$
    \H(r) = \{l$_0$:x$_{r0}$,$\ldots$,l$_n$:x$_{rn}$\} $\wedge$
    weak x$_{r0}$ = v$_0$ $\wedge$
    $\ldots$ $\wedge$
    weak x$_{rn}$ = v$_n$ }
}

\inferax{IStayEmpty}{\istay{$\emptyset$}{true}}
\inferax{IStayConst}{\istay{x=c}{weak x=c}}
\inferax{IStayObject}{\istay{x=\{l$_1$:v$_1$,$\ldots$,l$_n$:v$_n$\}}{weak x=\{l$_1$:v$_1$,$\ldots$,l$_n$:v$_n$\}}}
\inferax{IStayRef}{\istay{x=r}{required x=r}}

\inferrule{IStayOne}{
  \E$({\tt x}) = {\tt v}$ \andalso
  \E$_0 = $\E$ \setminus \{({\tt x},{\tt v})\}$ \andalso
  \istay{\E$_0$}{C$_0$} \andalso
  \istay{x=v}{C$_s$} \andalso
  {\tt C}$ = ${\tt C}$_0$ $\wedge$ {\tt C}$_s$
}{
  \istay{\E}{C}
}

\infrule[IStayHeap]{
  \mbox{
    \H$({\tt r}) = {\tt o}$ \andalso  
    \H$_0 = $\H$ \setminus \{({\tt r},{\tt o})\}$ \andalso
    \istay{\H$_0$}{C$_0$} \andalso
    ${\tt o} = $\{{\tt l$_0$:v$_0$,$\ldots$,l$_n$:v$_n$}\}}\\
  \mbox{
    \istay{x$_{r0}$=v$_0$}{C$_{r0}$} $\cdots$ \istay{x$_{rn}$=v$_n$}{C$_{rn}$}}
}{\mbox{
  \istay{\H}{C$_0$ $\wedge$
    \H(r) = \{l$_0$:x$_{r0}$,$\ldots$,l$_n$:x$_{rn}$\} $\wedge$
    C$_{r0}$ $\wedge$
    $\ldots$ $\wedge$
    C$_{rn}$ }
}}

\judgment{\inline{\E}{E}{\H}{\C}{\I}{e}{\E$'$}{e$_C$}{e$'$}}

``Inlining expression {\tt e} from the local environment {\tt E} turns
into expression {\tt e$'$}. To connect variables across method calls,
the constraint expression {\tt e$_C$} is returned.''

We use an inlining judgment to translate expressions into a
representation suitable for the solver. In particular, we translate
local variables into their names in the global environment and provide
a semantics for method calls inside constraints. Arguments to method
calls are constrained to be equal to the expression that generated
them. Inlining does not allow updates to the heap, so no new heap is
returned. We do allow assignments to locals in inlined methods,
however, so the global environment can change as a result of inlining.

%% TODO LATER:
%% \tf{SHOW: Any updates to the global environments would not affect
%%   existing variables. The rules only allow assigning to temps in
%%   methods that are inlined, so any changes to the global environment
%%   are not visible from any outside scope.}

\inferax{I-Const}{\inline{\E}{E}{\H}{\C}{\I}{c}{\E}{true}{c}}

\inferrule{I-Var}{
  {\tt E}({\tt x}) = {\tt x$_g$}
}{
  \inline{\E}{E}{\H}{\C}{\I}{x}{\E}{true}{x$_g$}
}

\inferrule{I-Value}{
  \inline{\E}{E}{\H}{\C}{\I}{e$_1$}{\E$_1$}{e$_{C_1}$}{e$'_1$} $\cdots$ \inline{\E}{E}{\H$_{n-1}$}{\C}{\I}{e$_n$}{\E$_n$}{e$_{C_n}$}{e$'_n$}
}{
  \inline{\E}{E}{\H}{\C}{\I}{\{l$_1$:e$_1$,$\ldots$,l$_n$:e$_n$\}}{\E$_n$}{e$_{C_1}$$\wedge\cdots\wedge$e$_{C_n}$}{\{l$_1$:e$'_1$,$\ldots$,l$_n$:e$'_n$\}}
}

\inferrule{I-Field}{
  \inline{\E}{E}{\H}{\C}{\I}{e}{\E$'$}{e$_C$}{e$'$} \andalso
  \evalBig{\E$'$}{E}{\H}{\C}{\I}{e}{\E$''$}{\H}{\C}{\I}{r}
}{
  \inline{\E}{E}{\H}{\C}{\I}{e.l}{\E$'$}{e$_C$$\wedge$e$'$=r}{\H(e$'$).l}
}

\inferrule{I-ValueField}{
  \inline{\E}{E}{\H}{\C}{\I}{e}{\E$'$}{e$_C$}{e$'$} \andalso
  \evalBig{\E$'$}{E}{\H}{\C}{\I}{e}{\E$''$}{\H}{\C}{\I}{\{l$_1$:v$_1$,$\ldots$,l$_n$:v$_n$\}}
}{
  \inline{\E}{E}{\H}{\C}{\I}{e.l}{\E$'$}{e$_C$}{e$'$.l}
}

\inferax{I-Ref}{\inline{\E}{E}{\H}{\C}{\I}{r}{\E}{true}{r}}

\inferrule{I-Op}{
  \inline{\E}{E}{\H}{\C}{\I}{e$_1$}{\E$'$}{e$_{C_a}$}{e$_a$} \andalso
  \inline{\E$'$}{E}{\H}{\C}{\I}{e$_2$}{\E$''$}{e$_{C_b}$}{e$_b$}
}{
  \inline{\E}{E}{\H}{\C}{\I}{e$_1$ $\oplus$ e$_2$}{\E$''$}{e$_{C_a}$$\wedge$e$_{C_b}$}{e$_a$ $\oplus$ e$_b$}
}

\inferrule{I-Compare}{
  \inline{\E}{E}{\H}{\C}{\I}{e$_1$}{\E$'$}{e$_{C_a}$}{e$_a$} \andalso
  \inline{\E$'$}{E}{\H}{\C}{\I}{e$_2$}{\E$''$}{e$_{C_b}$}{e$_b$}
}{
  \inline{\E}{E}{\H}{\C}{\I}{e$_1$ $\olessthan$ e$_2$}{\E$''$}{e$_{C_a}$$\wedge$e$_{C_b}$}{e$_a$ $\olessthan$ e$_b$}
}

\inferrule{I-Combine}{
  \inline{\E}{E}{\H}{\C}{\I}{e$_1$}{\E$'$}{e$_{C_a}$}{e$_a$} \andalso
  \inline{\E$'$}{E}{\H}{\C}{\I}{e$_2$}{\E$''$}{e$_{C_b}$}{e$_b$}
}{
  \inline{\E}{E}{\H}{\C}{\I}{e$_1$ $\ovee$ e$_2$}{\E$''$}{e$_{C_a}$$\wedge$e$_{C_b}$}{e$_a$ $\ovee$ e$_b$}
}

\inferrule{I-Identity}{
  \inline{\E}{E}{\H}{\C}{\I}{e$_1$}{\E$'$}{e$_{C_a}$}{e$_a$} \andalso
  \inline{\E$'$}{E}{\H}{\C}{\I}{e$_2$}{\E$''$}{e$_{C_b}$}{e$_b$}
}{
  \inline{\E}{E}{\H}{\C}{\I}{e$_1$ == e$_2$}{\E$''$}{e$_{C_a}$$\wedge$e$_{C_b}$}{e$_a$ == e$_b$}
}

\infrule[I-Call]{
  \mbox{
    \evalBig{\E}{E}{\H}{\C}{\I}{e}{\E$'$}{\H}{\C}{\I}{v} \andalso
    \lookup{v}{l}{s; return e}}\\
  \mbox{
    \enter{\E$'$}{E}{\H}{\C}{\I}{v}{\E$''$}{E$_m$}{\H}{\C}{\I}}\\
  \mbox{
    \stepHO{\E$''$}{E$_m$}{\H}{\C}{\I}{s}{\E$'''$}{E$'$}{\H}{\C}{\I} \andalso
    \evalBig{\E$'''$}{E$'$}{\H}{\C}{\I}{e}{\E$''''$}{\H}{\C}{\I}{v$_r$}}
}{\mbox{
  \inline{\E}{E}{\H}{\C}{\I}{e.l(e$_1$,$\ldots$,e$_n$)}{\E$''''$}{true}{v$_r$}
}}

Methods that have any statements at all can only be used in a one-way
manner. This is ensured by evaluating the return expression and using
only the value in the constraint. Because we are retranslating all
constraints on each semantic step, this return value will get updated
when its dependencies change, it just won't work in the other
direction.

When inlining a method with more than one statement, the statements
are simply executed.  In particular, this means that we eagerly choose
which branch of if-statements to inline and eagerly unroll loops.
Further, the {\sc I-Call} rule above and the {\sc I-MultiWayCall} rule
below ensure that methods being used in constraints have no side
effects.  This is accomplished by requiring the initial heap to remain
unchanged.

Similarly, methods used in constraints cannot declare nested
constraints; this is accomplished by requiring the initial sets of
ordinary and identity constraints to remain unchanged.

\infrule[I-MultiWayCall]{
  \mbox{
    \inline{\E}{E}{\H}{\C}{\I}{e$_0$}{\E$'$}{e$_{C_0}$}{e$_0'$} \andalso
    \evalBig{\E$'$}{E}{\H}{\C}{\I}{e$_0$}{\E$''$}{\H}{\C}{\I}{v}}\\
  \mbox{
    \lookup{v}{l}{return e}}\\
  \mbox{
    \enter{\E$''$}{E}{\H}{\C}{\I}{v}{\E$'''$}{E$_m$}{\H}{\C}{\I}}\\
  \mbox{
    \inline{\E$'''$}{E}{\H}{\C}{\I}{e$_1$}{\E$_1$}{e$_{C_1}$}{e$_1'$}
    $\cdots$
    \inline{\E$_{n-1}$}{E}{\H}{\C}{\I}{e$_n$}{\E$_n$}{e$_{C_n}$}{e$_n'$}}\\
  \mbox{
    {\tt E$_m$}({\tt self}) = {\tt x$_{g_{\tt self}}$} \andalso
    {\tt E$_m$}({\tt x$_1$}) = {\tt x$_{g_1}$} $\cdots$
    {\tt E$_m$}({\tt x$_n$}) = {\tt x$_{g_n}$}}\\
  \mbox{
    {\tt e$_C$} = ({\tt x$_{g_{\tt self}}$=e$_0'$ $\wedge$ x$_{g_1}$=e$_1'$ $\wedge$ $\cdots$ $\wedge$ x$_{g_n}$=e$_n'$}) \andalso
    \inline{\E$_n$}{E$_m$}{\H}{\C}{\I}{e}{\E$_n'$}{e$_{C_m}$}{e$'$}}
}{\mbox{
  \inline{\E}{E}{\H}{\C}{\I}{e$_0$.l(e$_1$,$\ldots$,e$_n$)}{\E$_n'$}{e$_C$$\wedge$e$_{C_m}$$\wedge$e$_{C_0}$$\wedge$e$_{C_1}$$\wedge\cdots\wedge$e$_{C_n}$}{e$'$}
}}

For methods that only return an expression, we inline the expression
and pass it to the solver. Note that we execute the argument
expressions and receiver for their value (potentially executing
through other methods), and also inline them, potentially executing
the same methods twice (once through \textsc{I-Call} and once through
\textsc{E-Call}.)  Although not ideal in terms of providing the
cleanest possible semantics, in practical terms this should 
not be a problem, because we prohibit
side-effects in these methods.


\judgment{\reinline{\E}{\H}{\I}{\C}{\E$'$}{C}}
\judgment{\reinline{\E}{\H}{\C}{\I}{\E$'$}{C}}

``Re-inlining the constraint store \C\ returns a
constraint {\tt C}''

``Re-inlining the constraint store \I\ returns a
constraint {\tt C}''

\inferax{I-ReinlineEmptyC}{\reinline{\E}{\H}{\I}{$\emptyset$}{\E}{true}}

\inferax{I-ReinlineEmptyI}{\reinline{\E}{\H}{\C}{$\emptyset$}{\E}{true}}

\inferrule{I-ReinlineC}{
  \C$_0$ = \C $\;\setminus$ $\{({\tt E}, \rho\, {\tt e})\}$ \andalso
  \reinline{\E}{\H}{\I}{\C$_0$}{\E$_0$}{C$_0$} \andalso
  \inline{\E$_0$}{E}{\H}{\C$_0$}{\I}{e}{\E$'$}{e$_{C_e}$}{e$'$}
}{
  \reinline{\E}{\H}{\I}{\C}{\E$'$}{C$_0$ $\wedge$ $\rho$ (e$'$ $\wedge$ e$_{C_e}$)}
}

\inferrule{I-ReinlineI}{
  \I$_0$ = \I $\;\setminus$ $\{({\tt E}, {\tt required\, e})\}$ \andalso
  \reinline{\E}{\H}{\C}{\I$_0$}{\E$_0$}{C$_0$} \andalso
  \inline{\E$_0$}{E}{\H}{\C}{\I$_0$}{e}{\E$'$}{e$_{C_e}$}{e$'$}
}{
  \reinline{\E}{\H}{\C}{\I}{\E$'$}{C$_0$ $\wedge$ required (e$'$ $\wedge$ e$_{C_e}$)}
}





\judgment{{\tt <\E|\H|\C|\I|I|C>} $\Longrightarrow$ {\tt <\E$'$|\H$'$>}}

This is a helper judgment for use in evaluating assignment statements
and identity constraints.  These statements are the only ones that can
cause the types of variables to change.  We handle this in two phases:
in a first phase, we propagate the new equality constraint through all
existing identity constraints in order to update other variables and
fields as needed.  In the second phase we solve all of the
non-identity constraints, plus any constraints {\tt C} created by
inlining identity constraints.

\infrule[TwoPhaseUpdate]{
  \mbox{
    \stay{\E}{C$_{E_s}$} \andalso
    \stay{\H}{C$_{{H}_s}$} \andalso
    \reinline{\E}{\H}{\C}{\I}{\E$_i$}{C$_i$} \andalso
  }\\
  \mbox{\solveH{\E$'$}{\H$'$}{(C$_i$ $\wedge$ C$_{E_s}$ $\wedge$ C$_0$ $\wedge$ C$_{{H}_s}$ $\wedge$ e$_1$=e$_2$)}}\\
  \mbox{
    \istay{\E$'$}{C$_{{E'}_s}$} \andalso
    \istay{\H$'$}{C$_{{H'}_s}$} \andalso
    \reinline{\E$'$}{\H$'$}{\I}{\C}{\E$_c$}{C} \andalso
    \tcHC{\E$_c$}{\H$'$}{C}}\\
  \mbox{
    \solveH{\E$''$}{\H$''$}{(C $\wedge$ C$_0$ $\wedge$ C$_{{E'}_s}$ $\wedge$ C$_{{H'}_s}$ $\wedge$ e$_1$=e$_2$)}}
}{\mbox{
  {\tt <\E|\H|\C|\I|e$_1$ == e$_2$|C$_0$>} $\Longrightarrow$ {\tt <\E$''$|\H$''$>}
}}


\judgment{\stepHO{\E}{E}{\H}{\C}{\I}{s}{\E$'$}{E$'$}{\H$'$}{\C$'$}{\I$'$}}

The stepping rules are refactored to work with the local environments
and the new constraint and identity-constraint stores.

\inferax{S-Skip}{
  \stepHO{\E}{E}{\H}{\C}{\I}{skip}{\E}{E}{\H}{\C}{\I}
}

\inferrule{S-Seq}{
  \stepHO{\E}{E}{\H}{\C}{\I}{s$_1$}{\E$'$}{E$'$}{\H$'$}{\C$'$}{\I$'$} \andalso
  \stepHO{\E$'$}{E$'$}{\H$'$}{\C$'$}{\I$'$}{s$_2$}{\E$''$}{E$''$}{\H$''$}{\C$''$}{\I$''$}
}{
  \stepHO{\E}{E}{\H}{\C}{\I}{s$_1$;s$_2$}{\E$''$}{E$''$}{\H$''$}{\C$''$}{\I$''$}
}

\inferrule{S-IfThen}{
  \evalBig{\E}{E}{\H}{\C}{\I}{e}{\E$'$}{\H$'$}{\C$'$}{\I$'$}{true} \andalso
  \stepHO{\E$'$}{E}{\H$'$}{\C$'$}{\I$'$}{s$_1$}{\E$''$}{E$'$}{\H$''$}{\C$''$}{\I$''$}
}{
  \stepHO{\E}{E}{\H}{\C}{\I}{if e then s$_1$ else s$_2$}{\E$''$}{E$'$}{\H$''$}{\C$''$}{\I$''$}%
}

\inferrule{S-IfElse}{
  \evalBig{\E}{E}{\H}{\C}{\I}{e}{\E$'$}{\H$'$}{\C$'$}{\I$'$}{false} \andalso
  \stepHO{\E$'$}{E}{\H$'$}{\C$'$}{\I$'$}{s$_2$}{\E$''$}{E$'$}{\H$''$}{\C$''$}{\I$''$}
}{
  \stepHO{\E}{E}{\H}{\C}{\I}{if e then s$_1$ else s$_2$}{\E$''$}{E$'$}{\H$''$}{\C$''$}{\I$''$}%
}

\infrule[S-WhileDo]{
  \mbox{
    \evalBig{\E}{E}{\H}{\C}{\I}{e}{\E$'$}{\H$'$}{\C$'$}{\I$'$}{true} \andalso
    \stepHO{\E$'$}{E}{\H$'$}{\C$'$}{\I$'$}{s}{\E$''$}{E$'$}{\H$''$}{\C$''$}{\I$''$}}\\
  \mbox{
    \stepHO{\E$''$}{E$'$}{\H$''$}{\C$''$}{\I$''$}{while e do s}{\E$'''$}{E$''$}{\H$'''$}{\C$'''$}{\I$'''$}}
}{\mbox{
  \stepHO{\E}{E}{\H}{\C}{\I}{while e do s}{\E$'''$}{E$''$}{\H$'''$}{\C$'''$}{\I$'''$}
}}

\inferrule{S-WhileSkip}{
  \evalBig{\E}{E}{\H}{\C}{\I}{e}{\E$'$}{\H$'$}{\C$'$}{\I$'$}{false}
}{
  \stepHO{\E}{E}{\H}{\C}{\I}{while e do s}{\E$'$}{E}{\H$'$}{\C$'$}{\I$'$}
}

\infrule[S-AsgnNewLocal]{
  \mbox{
    {\tt E}({\tt x})$\uparrow$ \andalso
    {\tt \E}({\tt x$_g$})$\uparrow$ \andalso
    ${\tt E}' = {\tt E} \bigcup \{({\tt x},{\tt x}_g)\}$}\\
  \mbox{
    \evalBig{\E}{E}{\H}{\C}{\I}{e}{\E$'$}{\H$'$}{\C$'$}{\I$'$}{v}}\\
  \mbox{
    {\tt <\E$'$|\H$'$|\C$'$|\I$'$|x$_g$==v|true>} $\Longrightarrow$ {\tt <\E$''$|\H$''$>}}
}{\mbox{
  \stepHO{\E}{E}{\H}{\C}{\I}{x := e}{\E$''$}{E$'$}{\H$''$}{\C$'$}{\I$'$}
}}

\infrule[S-AsgnLocal]{
  \mbox{
    ${\tt E}({\tt x}) = {\tt x}_g$}\\
  \mbox{
    \evalBig{\E}{E}{\H}{\C}{\I}{e}{\E$'$}{\H$'$}{\C$'$}{\I$'$}{v}}\\
  \mbox{
    {\tt <\E$'$|\H$'$|\C$'$|\I$'$|x$_g$==v|true>} $\Longrightarrow$ {\tt <\E$''$|\H$''$>}}
}{\mbox{
  \stepHO{\E}{E}{\H}{\C}{\I}{x := e}{\E$''$}{E}{\H$''$}{\C}{\I}
}}

\infrule[S-AsgnLValue]{
  \mbox{
    \evalBig{\E}{E}{\H}{\C}{\I}{e}{\E$'$}{\H$'$}{\C$'$}{\I$'$}{v} \andalso
    \inline{\E$'$}{E}{\H$'$}{\C$'$}{\I$'$}{e$_l$.l}{\E$''$}{e$_c$}{e$'$}} \\
  \mbox{
    {\tt <\E$''$|\H$'$|\C$'$|\I$'$|e$'$==v|e$_c$>} $\Longrightarrow$ {\tt <\E$'''$|\H$''$>}}
}{\mbox{
  \stepHO{\E}{E}{\H}{\C}{\I}{e$_l$.l := e}{\E$'''$}{E}{\H$''$}{\C$'$}{\I$'$}
}}

Note that the above rule can only be used in constraint-construction
mode if \H$ = $\H$'$\@.  This effectively restricts assignments in
constraints to work on values. Note also, that assignments in
constraints are just executed, but do not set up a required equality
constraint between the LHS and RHS\@.  (The current Babelsberg/Ruby
includes a ``backwards compatibility mode'' which sets up required
equality constraints when assignment occurs, but it's not clear
whether we should retain support for that.  See Appendix
\ref{sec:backwards-compatibility-mode}.)

We also do not allow the use of the following rules for {\tt once} and
{\tt always} in constraint-construction mode, because the inlining
rules disallow updating the constraint store and heap.  As discussed
previously (Section~\ref{sec:additional-constraint-restrictions}), in
a practical implementation we might want to support benign side
effects in methods that are invoked by constraint expressions,
including methods that themselves create new constraints; but this is
not modeled by this semantics.

\infrule[S-Once]{
  \mbox{
    {\tt C$_0$} = $\rho$ {\tt e} \andalso
    \inline{\E}{E}{\H}{\C}{\I}{e}{\E$'$}{e$_{C_e}$}{e$'$} \andalso
    {\tt C$_0'$} = $\rho$ ({\tt e$'$} $\wedge$ {\tt e$_{C_e}$}) \andalso
    \tcHC{\E$'$}{\H}{C$_0'$} \andalso
  }\\
  \mbox{
    \istay{\E$'$}{C$_{E_s}$} \andalso
    \istay{\H}{C$_{H_s}$} \andalso
    \reinline{\E$'$}{\H}{\I}{\C}{\E$''$}{C} \andalso
    \solveH{\E$'''$}{\H$'$}{$($C $\wedge$ C$_{E_s}$ $\wedge$ C$_{H_s}$ $\wedge$ C$_0'$$)$}}
}
{\mbox{\stepHO{\E}{E}{\H}{\C}{\I}{once C$_0$}{\E$'''$}{E}{\H$'$}{\C}{\I}}}

\infrule[S-Always]{
  \mbox{
    \stepHO{\E}{E}{\H}{\C}{\I}{once C$_0$}{\E$'$}{E}{\H$'$}{\C}{\I}}\\
  \mbox{
    \C$'$ = \C $\;\bigcup$ $\{({\tt E}, {\tt C}_0)\}$}
}
{\mbox{\stepHO{\E}{E}{\H}{\C}{\I}{always C$_0$}{\E$'$}{E}{\H$'$}{\C$'$}{\I}}}

\infrule[S-OnceIdentity]{
  \mbox{
    \evalBig{\E}{E}{\H}{\C}{\I}{e$_0$}{\E$_0$}{\H}{\C}{\I}{v} \andalso
    \evalBig{\E$_0$}{E}{\H}{\C}{\I}{e$_1$}{\E$_1$}{\H}{\C}{\I}{v}}\\
  \mbox{
    \inline{\E$_1$}{E}{\H}{\C}{\I}{e$_0$}{\E$_2$}{e$_{C_0}$}{e$_0'$} \andalso
    \inline{\E$_2$}{E}{\H}{\C}{\I}{e$_1$}{\E$_3$}{e$_{C_1}$}{e$_1'$}}\\
}{\mbox{1
    \stepHO{\E}{E}{\H}{\C}{\I}{once e$_0$ == e$_1$}{\E$_3$}{E}{\H}{\C}{\I}
}}

\infrule[S-AlwaysIdentity]{
  \mbox{
    \stepHO{\E}{E}{\H}{\C}{\I}{once e$_0$ == e$_1$}{\E$'$}{E}{\H$'$}{\C}{\I}}\\
  \mbox{
    \I$'$ = \I $\;\bigcup$ $\{(${\tt E}$,\,${\tt e}$_0$ {\tt ==} {\tt e}$_1)\}$}
}{\mbox{
    \stepHO{\E}{E}{\H}{\C}{\I}{always e$_0$ == e$_1$}{\E$'$}{E}{\H$'$}{\C}{\I$'$}
}}



%% ALAN todo:
%% Check the formal semantics for CLP, cc and tcc languages, which also
%% have a notion of a constraint store --- while we are doing something
%% different in the Babelsberg semantics, a reader who knows that
%% literature shouldn't come away thinking we are missing something
%% obvious, e.g. in terms of a better way to do something, falling into a
%% well-known pitfall, etc.

\newpage

\bibliographystyle{plain}
\bibliography{constraints}

\appendix
\newpage

\section{Appendix}

This appendix contains some observations that are potentially of
interest but that are mostly orthogonal to the task at hand of
providing a formal semantics for Babelsberg; and also descriptions of
some alternatives that were considered and later dropped, but that
seem promising enough that we might come back to them some day.  (In
other words, it's a holding pen for material we probably should just
have deleted but couldn't quite bear to.)

\subsection{Warnings and Debugging}
\label{sec:warnings-and-debugging}

We don't deal with warnings and debugging in the formal semantics for
Babelsberg.  In a practical, usable implementation, however, including
such support will be key.

First, all of the places where the formal semantics gets stuck to
model an error state should have explicit, understandable error
messages in a real implementation.  A few additional issues:

A non-required \verb|always| constraint might undo the result of an
assignment.

\begin{verbatim}
x := 0;
always strong x=10;
x := 5;
\end{verbatim}

\verb|x| starts out as 0, then is 10 to satisfy the \verb|always|
constraint, then gets set to 5, and then at the next time step, the
\verb|always| constraint causes it to be set back to 10.  We may want
to issue a warning in this case in a practical implementation.  (This
would be a simple case to check for.)

As described previously, in the interests of predictability, the
system will not create a new (standard) object to satisfy constraints.
There should be an appropriate error message in such a case if the
unsatisfied constraints are required, and a warning if they are soft.

\subsection{The Perturbation Model vs.\ the Refinement Model}
\label{sec:perturbation-vs-refinement}

A key issue in integrating constraints with imperative programming is
how to represent variables that change over time.  Earlier versions of
Kaleidoscope
\cite{freeman-benson-phd,freeman-benson-iccl-1992,freeman-benson-ecoop-1992}
used a \emph{refinement model}, in which ordinary variables were
represented as a stream of ``pellucid variables,'' each holding a
value at a different time.  This was in turn adapted from Lucid
\cite{ashcroft-cacm-1977}.  Later versions of Kaleidoscope
\cite{lopez-ecoop-1994} used a \emph{perturbation model}, in which
variables were represented conventionally.  Assignment perturbed the
value of a variable, and then the constraints took over and adjusted
the values of other variables so that they were re-satisfied if
necessary.  The formal semantics described here uses a variation of
the perturbation model, which we believe is cleaner.  (In particular,
rather than just perturbing the value of a variable by changing it, we
model assignment as a \verb|once| constraint between the variable and
the value of the expression on the right hand side of an assignment
statement.)

The refinement model does provide additional capabilities --- for
example syntax for referring to both the current and previous values
of a variable --- but at a cost, both conceptually for the programmer
and also for the language implementor.  At least for the common cases,
we believe the refinement model and our current model provide the same
answers.  

The refinement model makes essential use of read-only annotations on
variables to prevent the present from affecting the past.  So before
launching into an example, we digress to discuss read-only
annotations.

Intuitively, when choosing a best solution to a set of constraints,
constraints should not be allowed to affect the choice of values for
their read-only variables, i.e., information can flow out of the
read-only variables, but not into them.  There is a formal,
declarative definition of read-only annotations in
\cite{borning-lisp-symbolic-computation-1992}, which was in turn
adapted from that in the ALPS flat committed-choice logic language
\cite{maher-iclp-1987}.  A \emph{one-way constraint} can be
represented by annotating all but one of the constrained variables as
read-only.  One can also annotate an expression as read-only --- this
is syntactic sugar for annotating each of its variables as read-only.

Here's a simple example of the effect of a read-only annotation on a
variable.  First consider these constraints:

\begin{tabular}{rl}
required & $x=y$ \\
strong   & $x=3$ \\
weak     & $y=4$
\end{tabular}

This has the solution $x \mapsto 3, y \mapsto 3$, since the strong
constraint trumps the weak one.  

But suppose we add a read-only annotation:

\begin{tabular}{rl}
required & $x=y?$ \\
strong   & $x=3$ \\
weak     & $y=4$
\end{tabular}

Because $y$ is read-only in the $x=y?$ constraint, the solver can't
use that constraint in determining the value of $y$, and so the
solution this time is $x \mapsto 4, y \mapsto 4$.

Further, if we make the constraint $x=3$ be required:

\begin{tabular}{rl}
required & $x=y?$ \\
required & $x=3$ \\
weak     & $y=4$
\end{tabular}

the constraints now have no solution: on the one hand, any solution
must satisfy the required $x=3$ constraint, but on the other, we still
can't use $x=y?$ in determining a value for $y$.  (This is a so-called
``blocked'' set of constraints
\cite{borning-lisp-symbolic-computation-1992}.)

With those prelimaries out of the way, consider again the example in
Section~\ref{sec:babelsberg-primitivetypes} of unsatisfiable
constraints arising from the following program:

\begin{verbatim}
x := 5;
always x<=10;
x := x+15
\end{verbatim}

This behavior we modeled with the rules for Babelsberg/PrimitiveTypes
is the same as that of the refinement model, in which the program
would be equivalent to these constraints:

\begin{eqnarray*}
\mbox{required\ \ } x_0 & = & 5            \\
\forall t>0 \mbox{\ required\ \ } x_t & \leq & 10  \\
\mbox{weak\ \ } x_1 & = & x_0?             \\
\mbox{weak\ \ } x_2 & = & x_1?             \\
\mbox{required\ \ } x_2 & = & x_1? + 15
\end{eqnarray*}

The read-only annotations in the refinement model serve the same role
as do the evaluation rules in our current model.  In the example, in
our current semantics we model the final assignment \verb|x:=x+15| by
first evaluating \verb|x+15| in the old environment, and then adding a
required \verb|once| constraint that \verb|x| be equal to that value.
In the refinement model, we model the final assignment as $x_2 = x_1?
+ 15$, where the read-only annotation on $x_1$ accomplishes the same
thing, by preventing the solver from changing $x_1$ to satisfy this
constraint even though the other constraint that gives $x_1$ the value
of 5 is only weak.  Similarly, stay constraints in the current
semantics are modeled as weak constraints equating the variable with
its current value; in the refinment model these are weak constraints
relating the variables representing the current and previous versions,
with the previous version annotated as read-only.

\subsection{Issues with Using Once Constraints in the Semantics 
  for Assignments}

In all of the Babelsbergs, in the semantics we model an assignment by
evaluating the right-hand-side, creating a new variable if needed for
the left-hand-side, setting it to the new value, and then solving the
set of constraints that include a \verb|once| constraint that the
left-hand-side equal that value.  This lets assignment interact
correctly with other constraints in the constraint store, if any.  For
example, as a consequence of this semantics, the following program
will get stuck in the formal semantics, or raise an
exception in a practical implementation:

\begin{verbatim}
x := 0;
always x=10;
x := 5;
\end{verbatim}

However, there are some subtle issues in connection with this representation.

The analogous program still gets stuck, as one would want, in
Babelsberg/Records or when using value classes, since records and
instances of value classes are immutable.

\begin{verbatim}
x := {a:0};
always x.a=10;
x := {a:5};
\end{verbatim}


But once we get to Babelsberg/UID and Babelsberg/Objects with ordinary
objects, the situation changes.  Consider the analogous program in
Babelsberg/UID:
\begin{verbatim}
x := new {a:0};
always x.a=10;
x := new {a:5};
\end{verbatim}

Here, the right-hand-side of the last assignment is evaluated, the
reference to the new record is assigned to \verb|x|, and then the
system immediately changes its \verb|a| field to 10 --- it is only the
reference that is immutable, not the record on the heap.  This is
arguably weird, and so one of the design alternatives we considered
was to make the right hand side be read-only (recursively).  However,
this might be computationally expensive, and would also cause
difficulties when creating circular structures (see below).  So in the
current design, we just evaluate the right-hand-side to a value, but
don't do anything further about making it recursively read-only.

\subsubsection{Circular Structures}
\label{sec:circular-structures}

There is an issue regarding creating circular structures if we convert
assignment statements to \verb|once| constraints, with the value on
the right hand side recursively annotated as read only.  (The same
issue arose in an earlier version of the semantics in which we simply
turned an assignment into a \verb|once| constraint with the right hand
side annotated as read only, instead of first evaluating the right
hand side.)  Consider:
\begin{verbatim}
c := Cons.new(10,nil);
c.cdr := c;
\end{verbatim}

If the second assignment is converted to the constraint
\verb|once c.cdr == c?| this is unsatisfiable.  However, there is a
simple workaround, namely to replace such an assignment with a
\verb|once| identity constraint without the read-only annotation:
\begin{verbatim}
c := Cons.new(10,nil);
once c.cdr == c
\end{verbatim}

This would be a bit strange.  But this solution doesn't require any
changes to the formal semantics, and addresses all the previous
issues. 

{\bf Selectively Eliding the Read-Only Annotation.}
A variant that provides a more standard syntax is to make the entire
object on the right-hand side of the assignment be read-only, except
for a field that is assigned to (if any).  The constraints would be
exactly the same for all the examples except for the circular
structures one:

\begin{verbatim}
c := Cons.new(10,nil);
c.cdr := c;
\end{verbatim}

Here, \verb|c.car| and the reference to \verb|c| itself are read-only, but not
\verb|c.cdr|.

{\bf Distinguishing Ownership from Reference.}  Another alternative to
making the entire object on the right-hand side of the assignment
read-only is to distinguish ownership of parts from references to
other objects.  If an instance variable is for an ``owned'' part, then
that instance variable is made read-only; but otherwise not.  For
compatibility with the host language, the default would be that
instance variables do not refer to owned parts; rather, this must be
declared explicitly.  (It only needs to be declared if there are
constraints on the parts or subparts.)

\subsection{Alternatives to Structural Compatibility Checks}
\label{sec:structural-compatibility-check-alternatives}

Our model includes structural compatibility checks on constraints,
which prevent the solver from potentially generating new kinds of
records and other odd behaviors.  But since these ideas seem to keep
resurfacing, we outline some of our rejected alternatives in this
subsection, in case they end up being useful later.

We first consider records in Babelsberg/Records (immutable records, no UIDs).

If the solver will need to compare two records with different fields, we
need to extend the comparators to handle them.  For LPB, there are at
least two possibilities:
\begin{enumerate}
\item An equality constraint between records is either satisfied or
  it's not. For example consider a constraint \verb|p=q| when
  \verb|p={x:0, y:1}| and \verb|q={x:0, y:2}|, versus when
  \verb|q={a:1000}|.  In both cases the constraint is unsatisfied, and
  there is no reason to prefer one of these solutions over the other.

\item An equality constraint between records is unfolded into multiple
  constraints on the fields, each of which is satisfied or not.  (Such
  a constraint might be unsatisfied either because the values in one
  or more corresponding fields weren't equal, or were of different
  types, or because the corresponding field didn't even exist.)  In
  the above example, for the first case we have \verb|p.x=q.x|
  (satisfied) and \verb|p.y=q.y| (unsatisfied).  For the second case
  we have \verb|p.x=q.x| (unsatisfied) and \verb|p.y=q.y|
  (unsatisfied), and an additional constraint \verb|p.a=q.a|
  (unsatisfied).  The first solution would be preferred over the
  second under LPB.
\end{enumerate}

We also considered WSB for records.  The error for \verb|{x:5}={x:5}|
or \verb|{x:5}={x:6}| is clear enough (0 and 1 respectively).  But
what is the error for \verb|{x:5}={x:5,y:6}| or \verb|{x:5}={y:6}| or
\verb|{x:5}={x:"squid"}|?  We need to determine an error for
additional or missing fields, and for comparing different types and
this error needs to be such that it can be meaningfully compared with
the error for values of fields that hold numeric types.  Is the error
for \verb|{x:0}={x:1000000000}| less than for \verb|{x:0}={y:0}|?  We
could come up with a definition, but it's not clear it's that useful.

Another problem with having the solver handle records without
filtering out weird cases with structural compatibility checks is that
it may result in less predictability for programmers.  Consider:

\begin{verbatim}
p := {x:0};
q := {y:1};
always p=q;
\end{verbatim}

If this isn't rejected by a structural compatibility check, there
there are a number of plausible values that the solver could find for
{\tt p} and {\tt q} --- namely \verb|{x:0}|, \verb|{y:1}|, and
\verb|{x:0,y:1}| --- and unclear which is better.  The non-determinism
affects the structure of the result rather than just its value, and
would prevent achieving Goal \ref{item:deterministic} in our list in
Section \ref{sec:taming-identity-constraints}.

Yet another complication is that we would want a minimality condition
on records.  Suppose we let an \verb|always| constraint create a new
record if need be:
\begin{verbatim}
always p.x = 5;
\end{verbatim}
We want the solution to be \verb|p={x:5}|, but not e.g.
\verb|p={x:5, y:1000}|.

These issues, along the lack of clear use cases for having a solver
untamed by structural compatibility checks, led us to reject the
alternative outlined in this subsection.

For Babelsberg/Objects, the analog of this behavior would be to allow
the constraint solver generate new classes on the fly.  Particularly
since there would be multiple possible classes, as before, this seems
overly complex and unpredictable.

\subsection{Backwards Compatibility Mode for Methods}
\label{sec:backwards-compatibility-mode}

Methods called by constraints must be free of side effects.  (In a
practical language, benign side effects might be allowed, but we don't
model this in the formal semantics.)  This seems like a reasonable
restriction that we shouldn't attempt to change, although we probably
ought to make precise what constitutes a benign side effect.

In addition, if a method consists of more than just a single return
statement, it can only be used in the forward direction in
constraints.  This provides a straightforward way for both the
programmer and the implementor to understand how this should be
handled: if the method is used in the forward direction only, it can
be called in the ordinary way; if it is potentially used
multi-directionally, it is exploded.

If we are trying to add constraints to an existing language with a
substantial class library, this restriction does imply that many
methods, even though they are side-effect-free, can only be used in
the forward direction.  The Babelsberg/Ruby implementation provides a
``backwards compatibility mode'' that relaxes the restriction in an
attempt to make more of these existing methods be usable
multi-directionally.  (It won't allow an arbitrary method to be
successfully used in all modes in a constraint, but it increases the
chances that it will work.)  However, this mode is incompatible with
the new semantics for Babelsberg described in this memo.  We describe
it here for completeness; updating it to be compatible with the new
semantics (or simply removing it) is left for future work.

The Babelsberg/Ruby implementation uses two interpreter modes:
standard imperative execution mode and constraint construction mode.
When methods are evaluated in constraint construction mode, assignment
statements are converted to two-way equality constraints.  However,
we've now decided that methods called from a constraint expression
cannot themselves create other constraints, or call further methods
that do so (Section \ref{sec:additional-constraint-restrictions}),
which rules out this transformation in its current form.

For the record, though, if we relaxed this restriction, we might want
to support this conversion.  It still doesn't seem entirely clean, so
we should probably still omit it from the formal semantics.  Instead,
a practical implementation of Babelsberg might elect to include a
backwards compatibility mode for methods that does this
transformation, as well as other transformations that increase the
chances the method can be used in multi-directional constraints.  In
addition to converting assignments to constraints, the transformation
should create fresh variables each time a variable is used on the
left-hand-side of an assignment, and systematically use that new
variable thereafter (until another assignment).  First, here is an
example that just uses straight-line code.

\begin{verbatim}
def add_and_double(x,y,z)
  sum := 0;
  sum := sum+x;
  sum := sum+y;
  sum := sum+z;
  return 2*sum;
end;
\end{verbatim}

In backward compatibility mode this is rewritten as:
\begin{verbatim}
def add_and_double(x,y,z)
  sum_1 = 0;
  sum_2 = sum_1+x;
  sum_3 = sum_2+y;
  sum_4 = sum_3+z;
  return 2*sum_4;
end;
\end{verbatim}

Notice that a fresh variable is introduced for each assignment, but
not for the return statement.  Now, if we evaluate the constraint 
\verb|always 100 = add_and_double(10,15,n)|, the system should satisfy
the constraint by making \verb|n| be 25.

Conditionals, in which a variable might or might not be assigned to,
can be handled by using a required equality constraint to pass
through the old value for the other branch of the conditional.  For
example:

\newpage

\begin{verbatim}
def maybe_double(x)
  ans := x;
  if x<10
    then ans := 2*ans;
  end;
  return ans;
end;
\end{verbatim}

This becomes:

\begin{verbatim}
def maybe_double(x)
  ans_1 = x;
  if x<10
    then ans_2 = 2*ans_1;
    else ans_2 = ans_1;
  end;
  return ans_2;
end;
\end{verbatim}

Note though that the semantics of conditional statements are unchanged
--- we still eagerly evaluate the test and select the appropriate
branch (with no backtracking to the other branch).

If we also had a way to explode methods on demand to support
recursion, we might also be able to automatically convert methods with
loops to recursions that could then be used multi-directionally.  Here
is a version of \verb|sum| that might be automatically generated in
this fashion:

\begin{verbatim}
def sum
  helper(0,i,0,ans);
  return ans;
end;

def helper(old_i, new_i, old_ans, new_ans)
  if old_i<self.length
    then 
      temp_ans = old_ans + self[old_i];
      temp_i = old_i + 1;
      helper(temp_i,new_i,temp_ans,new_ans);
    else 
      new_ans = old_ans;
      new_i = old_i;
  end;
end;
\end{verbatim}

Here, \verb|helper| is produced automatically from the \verb|while| by
parameterizing the helper method with the old and new values of
\verb|i| and \verb|ans|, and converting the \verb|while| to an
\verb|if| with a recursive call.  (Again, we still eagerly evaluate
the test in the conditional, i.e., we evaluate the original
\verb|while| statement a definite number of times --- this method
won't let us try different array lengths to satisfy the constraints.)

As we've been noting, these transformations won't work in all cases.
As an extreme example, if we have an \verb|encrypt| method that takes
a plaintext message \verb|m| and a public key \verb|k|, we are unlikey
to be able to run the method backwards to find the original message
given the cyphertext.  Coming up with a good set of transformations, a
specification of when they will work, and whether the result covers a
useful set of cases, is an open problem.  Certainly they work for
methods consisting just of a sequence of assignment statements and a
final return statement, but this seems like not that useful a class of
methods to run backwards.

For the record, here are some more straightforward examples of when
the tranformation doesn't work.

Consider an absolute value method for integers:

\begin{verbatim}
def abs()
  if self>=0 then return self else return 0-self;
end;
\end{verbatim}

Suppose we use it in this program:

\begin{verbatim}
x := 5;
y := 5;
always y=x.abs();
always x<0;
\end{verbatim}

This works up until the final statement.  However, when we try to
satisfy \verb|x<0|, because of the semantics of the \verb|if|
statement (which eagerly evaluates the test, with no backtracking), we
incorrectly conclude the constraints are unsatisfiable.

Instead, the \verb|abs| method as written should only be used in the
forwards direction; if we want one that also works backwards, we
should write it as a disjunction that can be turned over to the
solver (or add support for Prolog-style backtracking --- see the next
section). 

Another case that doesn't work arises when the method has a
conditional, and the two branches return different types.  Even if
these are both value classes, this still doesn't explode correctly
using the current rules.  For example, suppose that in addition to
standard points, we have a class \verb|ZeroPoint| representing a point
with its x and y both 0 (which then avoids storing those fields).
Also suppose that both kinds of points implement an \verb|optimize|
method that returns a new object, which is the point represented in
the optimal way.  For ordinary points, the method is:

\begin{verbatim}
def optimize()
  if self.x=0 && self.y==0 then return ZeroPoint() else return self;
end;
\end{verbatim}

For \verb|ZeroPoint|, \verb|optimize| just returns \verb|self|.  

Now consider this code that uses the \verb|optimize| method:

\begin{verbatim}
p := Point(5,5);
q := Point(5,5);
always q = p.optimize();
\end{verbatim}

If we implement this by exploding the \verb|optimize| method, the
existing transformations don't work since \verb|q| might be either a
\verb|Point| or a \verb|ZeroPoint|.  We might be able to get it to
work in this case, but the transformations become more complex and
probably more fragile, and it's not clear that the complexity is worth
the price.

\subsection{Adding New Solvers and Extending the Solver Language}
\label{sec:adding-solvers}

In a practical Babelsberg implementation, it is useful to be able to
add new solvers, if necessary extending the solver language so that
new kinds of constraints can be encoded and sent to the new solvers.
This doesn't seem to create any issues for the formal semantics, so we
note this only in this appendix.

One useful extension of this sort would be useful to include a solver
for geometric constraints, either in 2 or 3 dimensions.

Some extensions that already exist in the practical languages are a
solver (and constructs in the solver language) for finite domain
constraints, as in the clpfd library for SWI Prolog.  Other additional
solvers are available for strings and local propagation constraints, as
solved for example by DeltaBlue.  Supporting the latter required adding a
construct for declaring local propagation constraints (including the
propagation methods) in the source language, and also adding local
propagation constraints to the solver language.

In Babelsberg/JS, when a DeltaBlue constraint is constructed, and an
equality constraint is encountered with two variables of the same type
on either side, when processing a constraint the interpreter doesn't
descend further, but instead asserts the equality on those
variables. So, for example, for two points the constraint expression
does not desugar into multiple constraints on the x and y
variables. This works even if the programmer does not specify a
propagation function, because DeltaBlue has a default propagation
method for equality (and some other relations, such as scaling.)

Finally, we may be able to have a solver language and solver to encode
Prolog-like goals, which would support backtracking and Prolog-style
programming within Babelsberg.  A key observation here is that this
would be in a separate set of constraints --- we wouldn't try to
change the basic Babelsberg semantics to allow backtracking with
Babelsberg \verb|if| statements.

\end{document}

%%% Local Variables:
%%% mode: latex
%%% TeX-master: t
%%% End:

% LocalWords:  Babelsberg Felgentreff Millstein Borning VPRI JS Javascript UID
% LocalWords:  PrimitiveTypes datatypes booleans implementers comparator LPB Z3
% LocalWords:  comparators WSB DeltaBlue iff unsatisfiable Disjunctions boolean
% LocalWords:  disjunctions disjunction datatype unevaluated Prolog unboxing '
% LocalWords:  metavariable unary satisfiable StayEmpty StayOne IfThen IfElse '
% LocalWords:  WhileDo WhileSkip PrimitiveType typechecking typecheck runtime '
% LocalWords:  Asgn uid Smalltalk desugar nondeterminism formalisms accessors
% LocalWords:  Accessor Dereference Const IdentityTrue IdentityFalse lookup 's
% LocalWords:  StayHeap AsgnNew OnceIdentity AlwaysIdentity Scala superclass 's
% LocalWords:  inlined multi directionally inlining RHS globals hoc There'd 's
